\chapter{Getting Started}\label{cha:start}

On first use, keep the sample files as they are and go through the steps in this chapter once: compile, confirm that a PDF is produced, and check that the fonts and pages display correctly. Then fill in your own thesis details and replace the titles and text chapter by chapter. Recompile after every change, so that if an error appears it is easy to tell which edit caused it.

The source files of this tutorial all live in \code{contents/}, and each chapter there corresponds to a chapter in the PDF. When you see a figure, citation or code layout you want, open the matching chapter, copy the example, and swap in your own text and data. Leave syntax you are unsure about as it is, adjust it step by step against the compiled result, and reuse it in other chapters once it works.

\section{Installing and Compiling}

Before compiling and writing, install a \LaTeX{} distribution. The template is compiled with \XeLaTeX{}; pick \href{https://www.tug.org/texlive/}{TeX Live} (MacTeX on macOS, the TeX Live installer on Windows) or the lighter \href{https://yihui.org/tinytex/}{TinyTeX}. The lines below are different ways to install one; run whichever fits your system. \code{brew} and \code{winget} must already be installed.

\begin{lstlisting}[language=bash]
brew install --cask mactex-no-gui   # macOS
winget install --id MiKTeX.MiKTeX   # Windows
sudo apt install texlive-full       # Debian / Ubuntu
\end{lstlisting}

The small program that converts the config file, \code{scripts/gen.lua}, is run by the \code{texlua} that comes with \TeX{}, so you do not need to install Python or anything else. References use the \code{citation-style-language} package; a full \TeX{} Live already includes it, and with TinyTeX or a partial install you add it yourself:

\begin{lstlisting}[language=bash]
tlmgr install citation-style-language
\end{lstlisting}

Once installed, switch to the project root in a terminal, which is the folder that contains \code{main.tex}, and run:

\begin{lstlisting}[language=bash]
texlua scripts/gen.lua
latexmk -xelatex main.tex
\end{lstlisting}

The first line reads \code{config.yml} and the school settings and writes \code{.tw-thesis-config.tex}. \LaTeX{} cannot read YAML directly, so this has to be repeated after every change to the config file. The second line is the actual compile; when it finishes, \code{main.pdf} appears in the same folder. Open the PDF and check that Chinese displays correctly. The fonts shipped with the project sit in \code{fonts/} and are loaded by file from \code{tw-thesis.cls}, so there is no need to install them on the system.

If \code{make} is installed, the two lines above shrink to the following commands, which run the jobs defined in the project's \code{Makefile}:

\begin{lstlisting}[language=bash]
make            # write main.pdf (draft, no watermark)
make watch      # recompile on every save
make release    # final version, with the school watermark
make clean      # delete main.pdf and the build leftovers
\end{lstlisting}

While writing, run \code{make watch} in the project root to start \textbf{watch mode}. The command keeps watching the files; every time you edit and save the thesis text or the settings, \code{latexmk} recompiles and updates the same \code{main.pdf}, so you never have to run the compile command by hand. Keep that terminal open while you work, and press \code{Ctrl+C} to stop watching.

To produce a version with a watermark, stop watch mode first and then run \code{make release}. It adds

\begin{lstlisting}[language=bash]
-usepretex='\def\twwatermark{1}'
\end{lstlisting}

and draws the watermark from the image named in the config file. It finishes after compiling and does not keep watching. Note that \code{make watch} and \code{make release} both write \code{main.pdf}; if you start \code{make watch} again later, the PDF is overwritten by the draft without the watermark.

\section{Project Structure}

After downloading or cloning the project, take a moment to learn what each folder is for, so you always know where to look when you want to change something. \tblref{tbl:layout} lists every file and directory in the project:

\begin{tlt}[notes={To change format rules edit \texttt{schools/}; to change text edit \texttt{contents/}.}]{Directories and files}{tbl:layout}{L{0.24\linewidth}L{0.66\linewidth}}
Path & Purpose \\
\midrule
\code{config.yml} & Your thesis details: title, author, department, date, keywords \\
\code{main.tex} & Binding order: front matter, chapters, references, appendices \\
\code{contents/} & The thesis text. \code{front/} holds the abstract and acknowledgements, \code{back/} the appendices and the \code{.bib} file \\
\code{schools/} & Format settings per school. \code{default.yml} is the only source of defaults \\
\code{scripts/} & The Lua script, run with \code{texlua}, that turns \code{config.yml} into \LaTeX{} macros \\
\code{tw-thesis.cls}, \code{tw-thesis.sty} & The template itself: the document class plus tools for figures, three-line tables and code style; independent of any one thesis and normally left alone \\
\code{figures/} & Thesis figures and sample figures \\
\code{assets/} & Template material such as school watermarks \\
\code{fonts/} & Fonts shipped with the template \\
\code{docs/} & The original PDFs of each school's thesis regulations \\
\end{tlt}

You will mostly edit \code{config.yml}, \code{main.tex} and \code{contents/}. Put thesis figures in \code{figures/}; look at \code{tw-thesis.sty} only when you need to adjust the shared style of figures, tables or code. The \code{\documentclass{tw-thesis}} line at the top of \code{main.tex} loads these tools as well; the content files themselves need no loading commands and can use \code{\fig}, \code{tlt} and the rest directly.

\section{Editing the Config File}

Once compiling works, open \code{config.yml} in the project root and replace the title, name, department and date with your own. The commonly used fields are listed below; fill in the English fields of the original file too, or one of the two languages will be missing:

\begin{lstlisting}[language=yaml]
title: "分散式系統的一致性研究"
title-en: "A Study on Consistency in Distributed Systems"

author: "王小明"
author-en: "Wang, Xiao-Ming"
id: "111753001"

advisor: "陳大文"
advisor-en: "Chen, Da-Wen"

college: "資訊學院"
institute: "資訊科學系"

degree: "碩士"
type: "論文"

date: "2027-01-15"       # date on the cover
oral-date: "2026-12-20"  # defense date printed on the verification letter

keywords: ["分散式系統", "一致性"]
keywords-en: ["Distributed Systems", "Consistency"]
\end{lstlisting}

Dates use \code{YYYY-MM-DD}. The cover converts them to the ROC year and an English month automatically, and the verification letter picks up the defense date as well.

\code{degree} (Master/Doctor) and \code{type} (thesis/technical report) are used on the cover and the verification letter. When you change them, check the matching \code{degree-en} and \code{type-en} too. After saving, run \code{texlua scripts/gen.lua} and compile again, look at the cover to see that the Chinese and English details agree, then check the date on the verification letter.

\section{Applying a School's Format}

Add \code{school:} to \code{config.yml}, or remove the leading \code{#} from the line already there, and fill in the school code. For example, to use National Chengchi University's settings:

\begin{lstlisting}[language=yaml]
school: nccu   # refers to schools/nccu.yml
\end{lstlisting}

The school codes provided by the project are listed in \tblref{tbl:schools}. A code must match the file name in \code{schools/}, without \code{.yml}. When no school is given, the template uses the default format and leaves the school name as placeholder text.

\begin{tlt}[notes={The original regulation PDFs are kept in \texttt{docs/}; file names carry the version date.}]{Built-in school codes}{tbl:schools}{L{0.1\linewidth}L{0.3\linewidth}L{0.1\linewidth}L{0.3\linewidth}}
Code & School & Code & School \\
\midrule
\code{ntu} & National Taiwan University & \code{ncu} & National Central University \\
\code{nthu} & National Tsing Hua University & \code{nsysu} & National Sun Yat-sen University \\
\code{nycu} & National Yang Ming Chiao Tung University & \code{nchu} & National Chung Hsing University \\
\code{ncku} & National Cheng Kung University & \code{ccu} & National Chung Cheng University \\
\code{nccu} & National Chengchi University & \code{ntou} & National Taiwan Ocean University \\
\code{ntpu} & National Taipei University & \code{nuk} & National University of Kaohsiung \\
\code{ntnu} & National Taiwan Normal University & \code{ncue} & National Changhua University of Education \\
\code{nknu} & National Kaohsiung Normal University & & \\
\end{tlt}

Settings are applied in order, with later ones overriding earlier ones:

\begin{lstlisting}
schools/default.yml  <-  schools/<school>.yml  <-  config.yml
\end{lstlisting}

That means you can override any format rule directly in \code{config.yml} without touching the school files. The config has several levels: taking \code{layout} as an example, \code{layout} contains \code{margin}, and \code{margin} contains the four margins. Overriding works item by item: only the items you actually write in \code{config.yml} are replaced, and anything you leave out keeps its default.

For example, the snippet below changes only the top margin and keeps the other three from the school settings:

\begin{lstlisting}[language=yaml]
layout:
  margin:
    top: 3.5
\end{lstlisting}

If \code{scripts/gen.lua} reports \code{unknown key}, check the field named in the message. In the example below \code{fnot-size} is misspelled and should be \code{font-size}:

\begin{lstlisting}
error: config.yml: unknown key `layout.fnot-size` (see schools/default.yml)
\end{lstlisting}

Every available key and its default is written in \code{schools/default.yml}, each with a comment giving the unit; that file is the complete list of settings.

\section{The Main File}

A printed thesis has a fixed binding order: cover, verification letter, table of contents, acknowledgements, abstract, chapters, references, appendices. \code{main.tex} is that order. It holds no thesis text; it only says which file goes where and where chapter numbering starts.

In other words, the text is written in the files under \code{contents/} and \code{main.tex} decides their order; to change the binding order or to add or remove a page, edit this file and leave the content alone. The whole file is short, and here it is from start to finish:

\begin{lstlisting}[language=latex]
\documentclass{tw-thesis}

\begin{document}

\thesisfront

% front matter
\begin{acknowledgement}
撰寫致謝時，請打開 \code{contents/front/acknowledgement.tex}，保留 \code{acknowledgement} 環境，再把這些說明換成自己的文字。模板會自動加入標題，內文依一般段落書寫即可，無需另外加上章節編號。

可以先回想研究過程中受到哪些協助，例如討論研究問題、取得資料，或協助檢查論文內容，再依序寫下想感謝的人與具體事情。完成後，請確認姓名、職稱與計畫名稱是否正確，也可以重新讀一次，看看是否有遺漏需要致謝的對象。

若不需要致謝頁，請到 \code{main.tex} 移除引用這個檔案的 \code{\input}，再重新編譯確認裝訂順序。若只是要更換標題，例如改為「謝詞」，則在 \code{config.yml} 的 \code{sections:} 下設定 \code{acknowledgement: "謝詞"} 即可。
\end{acknowledgement}

\begin{abstractzh}
撰寫中文摘要時，請先打開 \code{contents/front/abstract.tex}，找到 \code{abstractzh} 環境，將環境內的說明換成自己的內容。外層的 \code{\begin} 與 \code{\end} 請保留，模板會自動加入「摘要」標題，段落之間空一行即可。

內容可以先交代研究問題，再依序說明方法、主要結果與結論。整理結果時，盡量寫出實際觀察到的差異或數據，讓讀者知道這項研究得到什麼發現。完成初稿後，可以暫時不看正文，單獨讀一次摘要，確認研究的目的與結果是否交代清楚。

摘要下方的關鍵字由 \code{config.yml} 帶入，請在 \code{keywords} 填寫中文關鍵字。英文摘要則寫在同一個檔案的 \code{abstracten} 環境，並搭配 \code{keywords-en} 使用。兩種摘要填好後，請重新編譯，檢查分頁與關鍵字的位置，再比對中英文的內容是否一致。
\end{abstractzh}

\begin{abstracten}
Open \code{contents/front/abstract.tex} and find the \code{abstracten} environment. Replace the text inside it with your English abstract, keeping the \code{\begin} and \code{\end} lines in place. Leave a blank line between paragraphs; the template will add the heading.

Begin with the research question, then describe the method, main findings, and conclusion. Include specific results where possible. Once the draft is ready, compare it with the Chinese abstract to check that the findings and terminology agree.

To change the keywords below the abstract, edit \code{keywords-en} in \code{config.yml}. Compile the document again after updating both abstracts, and check the page breaks and the position of the keywords.
\end{abstracten}


% notation
\begin{denotation}
需要整理縮寫與符號時，請打開 \code{contents/front/denotation.tex}，保留 \code{denotation} 環境，再將下面的範例換成論文使用的項目。每個詞條以 \code{\item} 開頭，方括號內放縮寫或符號，後面填入完整名稱與定義。

填寫完成後，請對照正文檢查大小寫、數學符號與用詞是否一致，並移除沒有使用的項目。如果論文不需要符號列表，可以到 \code{main.tex} 移除對應的 \code{\input}，重新編譯後就不會出現這一頁。

\begin{description}
\item[PDF] 編譯後輸出的文件格式。
\item[YAML] 本模板設定檔使用的格式，副檔名為 \code{.yml}。
\item[BibTeX] 參考文獻資料庫所使用的格式。
\item[$E$] 能量。
\item[$m$] 質量。
\item[$c$] 光速。
\end{description}
\end{denotation}


% chapters
\chapter{開始使用}\label{cha:start}

初次使用時，請先保留範例檔案，按照本章的步驟編譯一次，確認能正常產生 PDF，並檢查中文字型與頁面是否顯示正確。完成後，再填入自己的論文資料，逐章替換標題與內文。建議每改完一個部分就重新編譯，若出現錯誤，也比較容易找出是哪一處修改造成的。

這份教學的原始檔都放在 \code{contents/}，各章的內容與 PDF 中的章節相互對應。看到想使用的圖表、引用或程式碼排版時，可以打開對應章節，複製範例，再換成自己的文字與資料即可。不熟悉的語法可以先保留原樣，對照編譯結果逐項調整，確認效果後再套用到其他章節。

\section{安裝與編譯}

在正式開始編譯與撰寫之前，需先安裝 \LaTeX{} 發行版，以及模板用來轉換設定檔的 Python 3。這份模板使用 \XeLaTeX{} 編譯，發行版請選擇 \href{https://www.tug.org/texlive/}{TeX Live}（macOS 可用 MacTeX，Windows 可用 TeX Live 安裝程式），或較輕量的 \href{https://yihui.org/tinytex/}{TinyTeX}。下面是不同的安裝方式，依自己的環境選一行執行即可；使用 \code{brew} 或 \code{winget} 前，電腦上須先有對應的套件管理工具。

\begin{lstlisting}[language=bash]
brew install --cask mactex-no-gui   # macOS
winget install --id MiKTeX.MiKTeX   # Windows
sudo apt install texlive-full       # Debian / Ubuntu
\end{lstlisting}

Python 3 的部分，macOS 與大多數 Linux 已內建，Windows 可到 \href{https://www.python.org/downloads/}{python.org} 下載。另外還需要 \code{PyYAML} 套件：

\begin{lstlisting}[language=bash]
python3 -m pip install pyyaml
\end{lstlisting}

安裝後，在終端機切換到專案根目錄，也就是放有 \code{main.tex} 的資料夾，執行下列指令：

\begin{lstlisting}[language=bash]
python3 scripts/gen.py
latexmk -xelatex main.tex
\end{lstlisting}

第一行讀取 \code{config.yml} 與學校設定，產生 \code{.tw-thesis-config.tex}；因為 \LaTeX{} 不能直接讀 YAML，所以每次修改設定檔後都要重新執行一次。第二行才是編譯，完成後會在同一個資料夾產生 \code{main.pdf}。打開 PDF，確認中文能正常顯示。專案附帶的字型放在 \code{fonts/}，編譯時由 \code{tw-thesis.cls} 直接依檔案載入，不必另外安裝到系統。

若已安裝 \code{make}，上面兩行可以縮短成下列指令，執行專案 \code{Makefile} 中定義的編譯工作：

\begin{lstlisting}[language=bash]
make            # 輸出 main.pdf（草稿，沒有浮水印）
make watch      # 存檔就重新編譯
make release    # 定稿，加上學校浮水印
make clean      # 刪掉 main.pdf 與編譯過程的暫存檔
\end{lstlisting}

撰寫論文時，可以在專案根目錄執行 \code{make watch}，啟動\textbf{監視模式}。指令執行後會持續監看檔案；每次修改並儲存論文內容或設定，\code{latexmk} 就會重新編譯，成功後更新同一份 \code{main.pdf}，不必每次手動執行編譯指令。使用期間請保持該終端機開啟；要停止監看時，切回終端機並按下 \code{Ctrl+C}。

需要輸出含浮水印的版本時，請先停止監視模式，再執行 \code{make release}。該指令會加入

\begin{lstlisting}[language=bash]
-usepretex='\def\twwatermark{1}'
\end{lstlisting}

依設定檔指定的圖片產生浮水印，編譯完成後便會結束，不會持續監看檔案。請注意，\code{make watch} 與 \code{make release} 都使用 \code{main.pdf} 作為輸出檔名；若之後重新啟動 \code{make watch}，PDF 會再次被不含浮水印的草稿版本覆寫。

\section{專案結構}

下載或複製專案後，可以先花時間認識各個資料夾的用途，之後要改任何內容就知道該打開哪個資料夾或是檔案。\tblref{tbl:layout} 是整個專案的檔案與目錄一覽：

\begin{tlt}[notes={要改格式規則請改 \texttt{schools/}，要改文字內容請改 \texttt{contents/}。}]{目錄與檔案}{tbl:layout}{L{0.24\linewidth}L{0.66\linewidth}}
路徑 & 用途 \\
\midrule
\code{config.yml} & 你的論文資料：題目、作者、系所、日期、關鍵字 \\
\code{main.tex} & 裝訂順序：前置、各章、參考文獻、附錄 \\
\code{contents/} & 論文內容。\code{front/} 是摘要與致謝，\code{back/} 是附錄與 \code{.bib} \\
\code{schools/} & 各校格式設定。\code{default.yml} 是唯一的預設值來源 \\
\code{scripts/} & 把 \code{config.yml} 轉成 \LaTeX{} 巨集的 Python 程式 \\
\code{tw-thesis.cls}、\code{tw-thesis.sty} & 模板本體：文件類別，以及圖片、三線表、程式碼樣式等工具；跟哪一本論文無關，一般不用改 \\
\code{figures/} & 論文圖片與範例圖片 \\
\code{assets/} & 學校浮水印等模板素材 \\
\code{fonts/} & 模板附帶的字型 \\
\code{docs/} & 各校論文格式規範的原始 PDF \\
\end{tlt}

基本上主要會改 \code{config.yml}、\code{main.tex} 和 \code{contents/}。論文圖片放進 \code{figures/}；需要調整圖片、表格或程式碼等共用樣式時再看 \code{tw-thesis.sty}。\code{main.tex} 開頭的 \code{\documentclass{tw-thesis}} 會一併載入這些工具，內容檔案本身不需要再寫任何載入指令，直接使用 \code{\fig}、\code{tlt} 等即可。

\section{修改設定檔}

確認可以成功編譯後，打開根目錄的設定檔 \code{config.yml}，先把題目、姓名、系所與日期換成自己的資料。下面列出常用欄位；原檔中的英文欄位也要一併填寫，否則將會出現某一語言缺失的情況：

\begin{lstlisting}[language=yaml]
title: "分散式系統的一致性研究"
title-en: "A Study on Consistency in Distributed Systems"

author: "王小明"
author-en: "Wang, Xiao-Ming"
id: "111753001"

advisor: "陳大文"
advisor-en: "Chen, Da-Wen"

college: "資訊學院"
institute: "資訊科學系"

degree: "碩士"
type: "論文"

date: "2027-01-15"       # 封面上的日期
oral-date: "2026-12-20"  # 印在審定書上的口試日期

keywords: ["分散式系統", "一致性"]
keywords-en: ["Distributed Systems", "Consistency"]
\end{lstlisting}

日期格式預設用 \code{YYYY-MM-DD}。封面會自動換成民國年與英文月份，審定書也會自動套上口試日期。

\code{degree}（碩士／博士）和 \code{type}（論文／技術報告）會用在封面和審定書。修改時，也要檢查對應的 \code{degree-en} 和 \code{type-en}。存檔後重新執行 \code{python3 scripts/gen.py} 與編譯，先看封面上的中英文資料是否一致，再檢查審定書的日期。

\section{套用學校格式}

在 \code{config.yml} 加入 \code{school:}，或移除原本那行開頭的 \code{#}，再填入學校代碼。例如使用政大的設定：

\begin{lstlisting}[language=yaml]
school: nccu   # 對應 schools/nccu.yml
\end{lstlisting}

專案提供的學校代碼列在\tblref{tbl:schools}。代碼必須和 \code{schools/} 裡的檔名相同，不含 \code{.yml}。未指定學校時，模板使用預設格式，校名也會保留為佔位文字。

\begin{tlt}[notes={各校規範的原始 PDF 收在 \texttt{docs/}，檔名含版本日期。}]{內建學校代碼}{tbl:schools}{L{0.1\linewidth}L{0.3\linewidth}L{0.1\linewidth}L{0.3\linewidth}}
代碼 & 學校 & 代碼 & 學校 \\
\midrule
\code{ntu} & 國立臺灣大學 & \code{ncu} & 國立中央大學 \\
\code{nthu} & 國立清華大學 & \code{nsysu} & 國立中山大學 \\
\code{nycu} & 國立陽明交通大學 & \code{nchu} & 國立中興大學 \\
\code{ncku} & 國立成功大學 & \code{ccu} & 國立中正大學 \\
\code{nccu} & 國立政治大學 & \code{ntou} & 國立臺灣海洋大學 \\
\code{ntpu} & 國立臺北大學 & \code{nuk} & 國立高雄大學 \\
\code{ntnu} & 國立臺灣師範大學 & \code{ncue} & 國立彰化師範大學 \\
\code{nknu} & 國立高雄師範大學 & & \\
\end{tlt}

設定檔的套用順序是後面設定覆蓋前面設定：

\begin{lstlisting}
schools/default.yml  <-  schools/<school>.yml  <-  config.yml
\end{lstlisting}

也就是說，你可以在 \code{config.yml} 裡直接覆蓋任何格式規則，不必去改學校設定檔。設定檔有好幾層，以 \code{layout} 為例，\code{layout} 底下有 \code{margin}，\code{margin} 底下才是上下左右四個邊界。覆蓋時是一項一項比對，僅有實際寫在 \code{config.yml} 的那一項才會被換掉，若無實際指定則維持預設值。

例如下方範例只修改上邊界，其餘三邊仍沿用學校設定：

\begin{lstlisting}[language=yaml]
layout:
  margin:
    top: 3.5
\end{lstlisting}

如果執行 \code{scripts/gen.py} 時出現 \code{unknown key}，先檢查錯誤訊息指出的欄位。例如下面的 \code{fnot-size} 拼錯了，應改成 \code{font-size}：

\begin{lstlisting}
error: config.yml: unknown key `layout.fnot-size` (see schools/default.yml)
\end{lstlisting}

所有可用的 key 與預設值都寫在 \code{schools/default.yml}，每一項都有註解說明單位，該份檔案就是完整的設定清單。

\section{主要檔案}

紙本論文有固定的裝訂順序：封面、審定書、目次、致謝、摘要、正文、參考文獻、附錄。\code{main.tex} 就是這份順序表。它不放任何論文內容，只寫明哪一個檔案排在哪個位置，以及從哪裡開始編章號。

換句話說，內容寫在 \code{contents/} 的各個檔案裡，\code{main.tex} 決定它們的先後；要調整裝訂順序、增刪某一頁，都是改這個檔案，不必去動內容。整份檔案很短，開頭到結尾就是下面這些：

\begin{lstlisting}[language=latex]
\documentclass{tw-thesis}

\begin{document}

\thesisfront

% front matter
\begin{acknowledgement}
撰寫致謝時，請打開 \code{contents/front/acknowledgement.tex}，保留 \code{acknowledgement} 環境，再把這些說明換成自己的文字。模板會自動加入標題，內文依一般段落書寫即可，無需另外加上章節編號。

可以先回想研究過程中受到哪些協助，例如討論研究問題、取得資料，或協助檢查論文內容，再依序寫下想感謝的人與具體事情。完成後，請確認姓名、職稱與計畫名稱是否正確，也可以重新讀一次，看看是否有遺漏需要致謝的對象。

若不需要致謝頁，請到 \code{main.tex} 移除引用這個檔案的 \code{\input}，再重新編譯確認裝訂順序。若只是要更換標題，例如改為「謝詞」，則在 \code{config.yml} 的 \code{sections:} 下設定 \code{acknowledgement: "謝詞"} 即可。
\end{acknowledgement}

\begin{abstractzh}
撰寫中文摘要時，請先打開 \code{contents/front/abstract.tex}，找到 \code{abstractzh} 環境，將環境內的說明換成自己的內容。外層的 \code{\begin} 與 \code{\end} 請保留，模板會自動加入「摘要」標題，段落之間空一行即可。

內容可以先交代研究問題，再依序說明方法、主要結果與結論。整理結果時，盡量寫出實際觀察到的差異或數據，讓讀者知道這項研究得到什麼發現。完成初稿後，可以暫時不看正文，單獨讀一次摘要，確認研究的目的與結果是否交代清楚。

摘要下方的關鍵字由 \code{config.yml} 帶入，請在 \code{keywords} 填寫中文關鍵字。英文摘要則寫在同一個檔案的 \code{abstracten} 環境，並搭配 \code{keywords-en} 使用。兩種摘要填好後，請重新編譯，檢查分頁與關鍵字的位置，再比對中英文的內容是否一致。
\end{abstractzh}

\begin{abstracten}
Open \code{contents/front/abstract.tex} and find the \code{abstracten} environment. Replace the text inside it with your English abstract, keeping the \code{\begin} and \code{\end} lines in place. Leave a blank line between paragraphs; the template will add the heading.

Begin with the research question, then describe the method, main findings, and conclusion. Include specific results where possible. Once the draft is ready, compare it with the Chinese abstract to check that the findings and terminology agree.

To change the keywords below the abstract, edit \code{keywords-en} in \code{config.yml}. Compile the document again after updating both abstracts, and check the page breaks and the position of the keywords.
\end{abstracten}


% notation
\begin{denotation}
需要整理縮寫與符號時，請打開 \code{contents/front/denotation.tex}，保留 \code{denotation} 環境，再將下面的範例換成論文使用的項目。每個詞條以 \code{\item} 開頭，方括號內放縮寫或符號，後面填入完整名稱與定義。

填寫完成後，請對照正文檢查大小寫、數學符號與用詞是否一致，並移除沒有使用的項目。如果論文不需要符號列表，可以到 \code{main.tex} 移除對應的 \code{\input}，重新編譯後就不會出現這一頁。

\begin{description}
\item[PDF] 編譯後輸出的文件格式。
\item[YAML] 本模板設定檔使用的格式，副檔名為 \code{.yml}。
\item[BibTeX] 參考文獻資料庫所使用的格式。
\item[$E$] 能量。
\item[$m$] 質量。
\item[$c$] 光速。
\end{description}
\end{denotation}


% chapters
\chapter{開始使用}\label{cha:start}

初次使用時，請先保留範例檔案，按照本章的步驟編譯一次，確認能正常產生 PDF，並檢查中文字型與頁面是否顯示正確。完成後，再填入自己的論文資料，逐章替換標題與內文。建議每改完一個部分就重新編譯，若出現錯誤，也比較容易找出是哪一處修改造成的。

這份教學的原始檔都放在 \code{contents/}，各章的內容與 PDF 中的章節相互對應。看到想使用的圖表、引用或程式碼排版時，可以打開對應章節，複製範例，再換成自己的文字與資料即可。不熟悉的語法可以先保留原樣，對照編譯結果逐項調整，確認效果後再套用到其他章節。

\section{安裝與編譯}

在正式開始編譯與撰寫之前，需先安裝 \LaTeX{} 發行版，以及模板用來轉換設定檔的 Python 3。這份模板使用 \XeLaTeX{} 編譯，發行版請選擇 \href{https://www.tug.org/texlive/}{TeX Live}（macOS 可用 MacTeX，Windows 可用 TeX Live 安裝程式），或較輕量的 \href{https://yihui.org/tinytex/}{TinyTeX}。下面是不同的安裝方式，依自己的環境選一行執行即可；使用 \code{brew} 或 \code{winget} 前，電腦上須先有對應的套件管理工具。

\begin{lstlisting}[language=bash]
brew install --cask mactex-no-gui   # macOS
winget install --id MiKTeX.MiKTeX   # Windows
sudo apt install texlive-full       # Debian / Ubuntu
\end{lstlisting}

Python 3 的部分，macOS 與大多數 Linux 已內建，Windows 可到 \href{https://www.python.org/downloads/}{python.org} 下載。另外還需要 \code{PyYAML} 套件：

\begin{lstlisting}[language=bash]
python3 -m pip install pyyaml
\end{lstlisting}

安裝後，在終端機切換到專案根目錄，也就是放有 \code{main.tex} 的資料夾，執行下列指令：

\begin{lstlisting}[language=bash]
python3 scripts/gen.py
latexmk -xelatex main.tex
\end{lstlisting}

第一行讀取 \code{config.yml} 與學校設定，產生 \code{.tw-thesis-config.tex}；因為 \LaTeX{} 不能直接讀 YAML，所以每次修改設定檔後都要重新執行一次。第二行才是編譯，完成後會在同一個資料夾產生 \code{main.pdf}。打開 PDF，確認中文能正常顯示。專案附帶的字型放在 \code{fonts/}，編譯時由 \code{tw-thesis.cls} 直接依檔案載入，不必另外安裝到系統。

若已安裝 \code{make}，上面兩行可以縮短成下列指令，執行專案 \code{Makefile} 中定義的編譯工作：

\begin{lstlisting}[language=bash]
make            # 輸出 main.pdf（草稿，沒有浮水印）
make watch      # 存檔就重新編譯
make release    # 定稿，加上學校浮水印
make clean      # 刪掉 main.pdf 與編譯過程的暫存檔
\end{lstlisting}

撰寫論文時，可以在專案根目錄執行 \code{make watch}，啟動\textbf{監視模式}。指令執行後會持續監看檔案；每次修改並儲存論文內容或設定，\code{latexmk} 就會重新編譯，成功後更新同一份 \code{main.pdf}，不必每次手動執行編譯指令。使用期間請保持該終端機開啟；要停止監看時，切回終端機並按下 \code{Ctrl+C}。

需要輸出含浮水印的版本時，請先停止監視模式，再執行 \code{make release}。該指令會加入

\begin{lstlisting}[language=bash]
-usepretex='\def\twwatermark{1}'
\end{lstlisting}

依設定檔指定的圖片產生浮水印，編譯完成後便會結束，不會持續監看檔案。請注意，\code{make watch} 與 \code{make release} 都使用 \code{main.pdf} 作為輸出檔名；若之後重新啟動 \code{make watch}，PDF 會再次被不含浮水印的草稿版本覆寫。

\section{專案結構}

下載或複製專案後，可以先花時間認識各個資料夾的用途，之後要改任何內容就知道該打開哪個資料夾或是檔案。\tblref{tbl:layout} 是整個專案的檔案與目錄一覽：

\begin{tlt}[notes={要改格式規則請改 \texttt{schools/}，要改文字內容請改 \texttt{contents/}。}]{目錄與檔案}{tbl:layout}{L{0.24\linewidth}L{0.66\linewidth}}
路徑 & 用途 \\
\midrule
\code{config.yml} & 你的論文資料：題目、作者、系所、日期、關鍵字 \\
\code{main.tex} & 裝訂順序：前置、各章、參考文獻、附錄 \\
\code{contents/} & 論文內容。\code{front/} 是摘要與致謝，\code{back/} 是附錄與 \code{.bib} \\
\code{schools/} & 各校格式設定。\code{default.yml} 是唯一的預設值來源 \\
\code{scripts/} & 把 \code{config.yml} 轉成 \LaTeX{} 巨集的 Python 程式 \\
\code{tw-thesis.cls}、\code{tw-thesis.sty} & 模板本體：文件類別，以及圖片、三線表、程式碼樣式等工具；跟哪一本論文無關，一般不用改 \\
\code{figures/} & 論文圖片與範例圖片 \\
\code{assets/} & 學校浮水印等模板素材 \\
\code{fonts/} & 模板附帶的字型 \\
\code{docs/} & 各校論文格式規範的原始 PDF \\
\end{tlt}

基本上主要會改 \code{config.yml}、\code{main.tex} 和 \code{contents/}。論文圖片放進 \code{figures/}；需要調整圖片、表格或程式碼等共用樣式時再看 \code{tw-thesis.sty}。\code{main.tex} 開頭的 \code{\documentclass{tw-thesis}} 會一併載入這些工具，內容檔案本身不需要再寫任何載入指令，直接使用 \code{\fig}、\code{tlt} 等即可。

\section{修改設定檔}

確認可以成功編譯後，打開根目錄的設定檔 \code{config.yml}，先把題目、姓名、系所與日期換成自己的資料。下面列出常用欄位；原檔中的英文欄位也要一併填寫，否則將會出現某一語言缺失的情況：

\begin{lstlisting}[language=yaml]
title: "分散式系統的一致性研究"
title-en: "A Study on Consistency in Distributed Systems"

author: "王小明"
author-en: "Wang, Xiao-Ming"
id: "111753001"

advisor: "陳大文"
advisor-en: "Chen, Da-Wen"

college: "資訊學院"
institute: "資訊科學系"

degree: "碩士"
type: "論文"

date: "2027-01-15"       # 封面上的日期
oral-date: "2026-12-20"  # 印在審定書上的口試日期

keywords: ["分散式系統", "一致性"]
keywords-en: ["Distributed Systems", "Consistency"]
\end{lstlisting}

日期格式預設用 \code{YYYY-MM-DD}。封面會自動換成民國年與英文月份，審定書也會自動套上口試日期。

\code{degree}（碩士／博士）和 \code{type}（論文／技術報告）會用在封面和審定書。修改時，也要檢查對應的 \code{degree-en} 和 \code{type-en}。存檔後重新執行 \code{python3 scripts/gen.py} 與編譯，先看封面上的中英文資料是否一致，再檢查審定書的日期。

\section{套用學校格式}

在 \code{config.yml} 加入 \code{school:}，或移除原本那行開頭的 \code{#}，再填入學校代碼。例如使用政大的設定：

\begin{lstlisting}[language=yaml]
school: nccu   # 對應 schools/nccu.yml
\end{lstlisting}

專案提供的學校代碼列在\tblref{tbl:schools}。代碼必須和 \code{schools/} 裡的檔名相同，不含 \code{.yml}。未指定學校時，模板使用預設格式，校名也會保留為佔位文字。

\begin{tlt}[notes={各校規範的原始 PDF 收在 \texttt{docs/}，檔名含版本日期。}]{內建學校代碼}{tbl:schools}{L{0.1\linewidth}L{0.3\linewidth}L{0.1\linewidth}L{0.3\linewidth}}
代碼 & 學校 & 代碼 & 學校 \\
\midrule
\code{ntu} & 國立臺灣大學 & \code{ncu} & 國立中央大學 \\
\code{nthu} & 國立清華大學 & \code{nsysu} & 國立中山大學 \\
\code{nycu} & 國立陽明交通大學 & \code{nchu} & 國立中興大學 \\
\code{ncku} & 國立成功大學 & \code{ccu} & 國立中正大學 \\
\code{nccu} & 國立政治大學 & \code{ntou} & 國立臺灣海洋大學 \\
\code{ntpu} & 國立臺北大學 & \code{nuk} & 國立高雄大學 \\
\code{ntnu} & 國立臺灣師範大學 & \code{ncue} & 國立彰化師範大學 \\
\code{nknu} & 國立高雄師範大學 & & \\
\end{tlt}

設定檔的套用順序是後面設定覆蓋前面設定：

\begin{lstlisting}
schools/default.yml  <-  schools/<school>.yml  <-  config.yml
\end{lstlisting}

也就是說，你可以在 \code{config.yml} 裡直接覆蓋任何格式規則，不必去改學校設定檔。設定檔有好幾層，以 \code{layout} 為例，\code{layout} 底下有 \code{margin}，\code{margin} 底下才是上下左右四個邊界。覆蓋時是一項一項比對，僅有實際寫在 \code{config.yml} 的那一項才會被換掉，若無實際指定則維持預設值。

例如下方範例只修改上邊界，其餘三邊仍沿用學校設定：

\begin{lstlisting}[language=yaml]
layout:
  margin:
    top: 3.5
\end{lstlisting}

如果執行 \code{scripts/gen.py} 時出現 \code{unknown key}，先檢查錯誤訊息指出的欄位。例如下面的 \code{fnot-size} 拼錯了，應改成 \code{font-size}：

\begin{lstlisting}
error: config.yml: unknown key `layout.fnot-size` (see schools/default.yml)
\end{lstlisting}

所有可用的 key 與預設值都寫在 \code{schools/default.yml}，每一項都有註解說明單位，該份檔案就是完整的設定清單。

\section{主要檔案}

紙本論文有固定的裝訂順序：封面、審定書、目次、致謝、摘要、正文、參考文獻、附錄。\code{main.tex} 就是這份順序表。它不放任何論文內容，只寫明哪一個檔案排在哪個位置，以及從哪裡開始編章號。

換句話說，內容寫在 \code{contents/} 的各個檔案裡，\code{main.tex} 決定它們的先後；要調整裝訂順序、增刪某一頁，都是改這個檔案，不必去動內容。整份檔案很短，開頭到結尾就是下面這些：

\begin{lstlisting}[language=latex]
\documentclass{tw-thesis}

\begin{document}

\thesisfront

% front matter
\begin{acknowledgement}
撰寫致謝時，請打開 \code{contents/front/acknowledgement.tex}，保留 \code{acknowledgement} 環境，再把這些說明換成自己的文字。模板會自動加入標題，內文依一般段落書寫即可，無需另外加上章節編號。

可以先回想研究過程中受到哪些協助，例如討論研究問題、取得資料，或協助檢查論文內容，再依序寫下想感謝的人與具體事情。完成後，請確認姓名、職稱與計畫名稱是否正確，也可以重新讀一次，看看是否有遺漏需要致謝的對象。

若不需要致謝頁，請到 \code{main.tex} 移除引用這個檔案的 \code{\input}，再重新編譯確認裝訂順序。若只是要更換標題，例如改為「謝詞」，則在 \code{config.yml} 的 \code{sections:} 下設定 \code{acknowledgement: "謝詞"} 即可。
\end{acknowledgement}

\begin{abstractzh}
撰寫中文摘要時，請先打開 \code{contents/front/abstract.tex}，找到 \code{abstractzh} 環境，將環境內的說明換成自己的內容。外層的 \code{\begin} 與 \code{\end} 請保留，模板會自動加入「摘要」標題，段落之間空一行即可。

內容可以先交代研究問題，再依序說明方法、主要結果與結論。整理結果時，盡量寫出實際觀察到的差異或數據，讓讀者知道這項研究得到什麼發現。完成初稿後，可以暫時不看正文，單獨讀一次摘要，確認研究的目的與結果是否交代清楚。

摘要下方的關鍵字由 \code{config.yml} 帶入，請在 \code{keywords} 填寫中文關鍵字。英文摘要則寫在同一個檔案的 \code{abstracten} 環境，並搭配 \code{keywords-en} 使用。兩種摘要填好後，請重新編譯，檢查分頁與關鍵字的位置，再比對中英文的內容是否一致。
\end{abstractzh}

\begin{abstracten}
Open \code{contents/front/abstract.tex} and find the \code{abstracten} environment. Replace the text inside it with your English abstract, keeping the \code{\begin} and \code{\end} lines in place. Leave a blank line between paragraphs; the template will add the heading.

Begin with the research question, then describe the method, main findings, and conclusion. Include specific results where possible. Once the draft is ready, compare it with the Chinese abstract to check that the findings and terminology agree.

To change the keywords below the abstract, edit \code{keywords-en} in \code{config.yml}. Compile the document again after updating both abstracts, and check the page breaks and the position of the keywords.
\end{abstracten}


% notation
\begin{denotation}
需要整理縮寫與符號時，請打開 \code{contents/front/denotation.tex}，保留 \code{denotation} 環境，再將下面的範例換成論文使用的項目。每個詞條以 \code{\item} 開頭，方括號內放縮寫或符號，後面填入完整名稱與定義。

填寫完成後，請對照正文檢查大小寫、數學符號與用詞是否一致，並移除沒有使用的項目。如果論文不需要符號列表，可以到 \code{main.tex} 移除對應的 \code{\input}，重新編譯後就不會出現這一頁。

\begin{description}
\item[PDF] 編譯後輸出的文件格式。
\item[YAML] 本模板設定檔使用的格式，副檔名為 \code{.yml}。
\item[BibTeX] 參考文獻資料庫所使用的格式。
\item[$E$] 能量。
\item[$m$] 質量。
\item[$c$] 光速。
\end{description}
\end{denotation}


% chapters
\chapter{開始使用}\label{cha:start}

初次使用時，請先保留範例檔案，按照本章的步驟編譯一次，確認能正常產生 PDF，並檢查中文字型與頁面是否顯示正確。完成後，再填入自己的論文資料，逐章替換標題與內文。建議每改完一個部分就重新編譯，若出現錯誤，也比較容易找出是哪一處修改造成的。

這份教學的原始檔都放在 \code{contents/}，各章的內容與 PDF 中的章節相互對應。看到想使用的圖表、引用或程式碼排版時，可以打開對應章節，複製範例，再換成自己的文字與資料即可。不熟悉的語法可以先保留原樣，對照編譯結果逐項調整，確認效果後再套用到其他章節。

\section{安裝與編譯}

在正式開始編譯與撰寫之前，需先安裝 \LaTeX{} 發行版，以及模板用來轉換設定檔的 Python 3。這份模板使用 \XeLaTeX{} 編譯，發行版請選擇 \href{https://www.tug.org/texlive/}{TeX Live}（macOS 可用 MacTeX，Windows 可用 TeX Live 安裝程式），或較輕量的 \href{https://yihui.org/tinytex/}{TinyTeX}。下面是不同的安裝方式，依自己的環境選一行執行即可；使用 \code{brew} 或 \code{winget} 前，電腦上須先有對應的套件管理工具。

\begin{lstlisting}[language=bash]
brew install --cask mactex-no-gui   # macOS
winget install --id MiKTeX.MiKTeX   # Windows
sudo apt install texlive-full       # Debian / Ubuntu
\end{lstlisting}

Python 3 的部分，macOS 與大多數 Linux 已內建，Windows 可到 \href{https://www.python.org/downloads/}{python.org} 下載。另外還需要 \code{PyYAML} 套件：

\begin{lstlisting}[language=bash]
python3 -m pip install pyyaml
\end{lstlisting}

安裝後，在終端機切換到專案根目錄，也就是放有 \code{main.tex} 的資料夾，執行下列指令：

\begin{lstlisting}[language=bash]
python3 scripts/gen.py
latexmk -xelatex main.tex
\end{lstlisting}

第一行讀取 \code{config.yml} 與學校設定，產生 \code{.tw-thesis-config.tex}；因為 \LaTeX{} 不能直接讀 YAML，所以每次修改設定檔後都要重新執行一次。第二行才是編譯，完成後會在同一個資料夾產生 \code{main.pdf}。打開 PDF，確認中文能正常顯示。專案附帶的字型放在 \code{fonts/}，編譯時由 \code{tw-thesis.cls} 直接依檔案載入，不必另外安裝到系統。

若已安裝 \code{make}，上面兩行可以縮短成下列指令，執行專案 \code{Makefile} 中定義的編譯工作：

\begin{lstlisting}[language=bash]
make            # 輸出 main.pdf（草稿，沒有浮水印）
make watch      # 存檔就重新編譯
make release    # 定稿，加上學校浮水印
make clean      # 刪掉 main.pdf 與編譯過程的暫存檔
\end{lstlisting}

撰寫論文時，可以在專案根目錄執行 \code{make watch}，啟動\textbf{監視模式}。指令執行後會持續監看檔案；每次修改並儲存論文內容或設定，\code{latexmk} 就會重新編譯，成功後更新同一份 \code{main.pdf}，不必每次手動執行編譯指令。使用期間請保持該終端機開啟；要停止監看時，切回終端機並按下 \code{Ctrl+C}。

需要輸出含浮水印的版本時，請先停止監視模式，再執行 \code{make release}。該指令會加入

\begin{lstlisting}[language=bash]
-usepretex='\def\twwatermark{1}'
\end{lstlisting}

依設定檔指定的圖片產生浮水印，編譯完成後便會結束，不會持續監看檔案。請注意，\code{make watch} 與 \code{make release} 都使用 \code{main.pdf} 作為輸出檔名；若之後重新啟動 \code{make watch}，PDF 會再次被不含浮水印的草稿版本覆寫。

\section{專案結構}

下載或複製專案後，可以先花時間認識各個資料夾的用途，之後要改任何內容就知道該打開哪個資料夾或是檔案。\tblref{tbl:layout} 是整個專案的檔案與目錄一覽：

\begin{tlt}[notes={要改格式規則請改 \texttt{schools/}，要改文字內容請改 \texttt{contents/}。}]{目錄與檔案}{tbl:layout}{L{0.24\linewidth}L{0.66\linewidth}}
路徑 & 用途 \\
\midrule
\code{config.yml} & 你的論文資料：題目、作者、系所、日期、關鍵字 \\
\code{main.tex} & 裝訂順序：前置、各章、參考文獻、附錄 \\
\code{contents/} & 論文內容。\code{front/} 是摘要與致謝，\code{back/} 是附錄與 \code{.bib} \\
\code{schools/} & 各校格式設定。\code{default.yml} 是唯一的預設值來源 \\
\code{scripts/} & 把 \code{config.yml} 轉成 \LaTeX{} 巨集的 Python 程式 \\
\code{tw-thesis.cls}、\code{tw-thesis.sty} & 模板本體：文件類別，以及圖片、三線表、程式碼樣式等工具；跟哪一本論文無關，一般不用改 \\
\code{figures/} & 論文圖片與範例圖片 \\
\code{assets/} & 學校浮水印等模板素材 \\
\code{fonts/} & 模板附帶的字型 \\
\code{docs/} & 各校論文格式規範的原始 PDF \\
\end{tlt}

基本上主要會改 \code{config.yml}、\code{main.tex} 和 \code{contents/}。論文圖片放進 \code{figures/}；需要調整圖片、表格或程式碼等共用樣式時再看 \code{tw-thesis.sty}。\code{main.tex} 開頭的 \code{\documentclass{tw-thesis}} 會一併載入這些工具，內容檔案本身不需要再寫任何載入指令，直接使用 \code{\fig}、\code{tlt} 等即可。

\section{修改設定檔}

確認可以成功編譯後，打開根目錄的設定檔 \code{config.yml}，先把題目、姓名、系所與日期換成自己的資料。下面列出常用欄位；原檔中的英文欄位也要一併填寫，否則將會出現某一語言缺失的情況：

\begin{lstlisting}[language=yaml]
title: "分散式系統的一致性研究"
title-en: "A Study on Consistency in Distributed Systems"

author: "王小明"
author-en: "Wang, Xiao-Ming"
id: "111753001"

advisor: "陳大文"
advisor-en: "Chen, Da-Wen"

college: "資訊學院"
institute: "資訊科學系"

degree: "碩士"
type: "論文"

date: "2027-01-15"       # 封面上的日期
oral-date: "2026-12-20"  # 印在審定書上的口試日期

keywords: ["分散式系統", "一致性"]
keywords-en: ["Distributed Systems", "Consistency"]
\end{lstlisting}

日期格式預設用 \code{YYYY-MM-DD}。封面會自動換成民國年與英文月份，審定書也會自動套上口試日期。

\code{degree}（碩士／博士）和 \code{type}（論文／技術報告）會用在封面和審定書。修改時，也要檢查對應的 \code{degree-en} 和 \code{type-en}。存檔後重新執行 \code{python3 scripts/gen.py} 與編譯，先看封面上的中英文資料是否一致，再檢查審定書的日期。

\section{套用學校格式}

在 \code{config.yml} 加入 \code{school:}，或移除原本那行開頭的 \code{#}，再填入學校代碼。例如使用政大的設定：

\begin{lstlisting}[language=yaml]
school: nccu   # 對應 schools/nccu.yml
\end{lstlisting}

專案提供的學校代碼列在\tblref{tbl:schools}。代碼必須和 \code{schools/} 裡的檔名相同，不含 \code{.yml}。未指定學校時，模板使用預設格式，校名也會保留為佔位文字。

\begin{tlt}[notes={各校規範的原始 PDF 收在 \texttt{docs/}，檔名含版本日期。}]{內建學校代碼}{tbl:schools}{L{0.1\linewidth}L{0.3\linewidth}L{0.1\linewidth}L{0.3\linewidth}}
代碼 & 學校 & 代碼 & 學校 \\
\midrule
\code{ntu} & 國立臺灣大學 & \code{ncu} & 國立中央大學 \\
\code{nthu} & 國立清華大學 & \code{nsysu} & 國立中山大學 \\
\code{nycu} & 國立陽明交通大學 & \code{nchu} & 國立中興大學 \\
\code{ncku} & 國立成功大學 & \code{ccu} & 國立中正大學 \\
\code{nccu} & 國立政治大學 & \code{ntou} & 國立臺灣海洋大學 \\
\code{ntpu} & 國立臺北大學 & \code{nuk} & 國立高雄大學 \\
\code{ntnu} & 國立臺灣師範大學 & \code{ncue} & 國立彰化師範大學 \\
\code{nknu} & 國立高雄師範大學 & & \\
\end{tlt}

設定檔的套用順序是後面設定覆蓋前面設定：

\begin{lstlisting}
schools/default.yml  <-  schools/<school>.yml  <-  config.yml
\end{lstlisting}

也就是說，你可以在 \code{config.yml} 裡直接覆蓋任何格式規則，不必去改學校設定檔。設定檔有好幾層，以 \code{layout} 為例，\code{layout} 底下有 \code{margin}，\code{margin} 底下才是上下左右四個邊界。覆蓋時是一項一項比對，僅有實際寫在 \code{config.yml} 的那一項才會被換掉，若無實際指定則維持預設值。

例如下方範例只修改上邊界，其餘三邊仍沿用學校設定：

\begin{lstlisting}[language=yaml]
layout:
  margin:
    top: 3.5
\end{lstlisting}

如果執行 \code{scripts/gen.py} 時出現 \code{unknown key}，先檢查錯誤訊息指出的欄位。例如下面的 \code{fnot-size} 拼錯了，應改成 \code{font-size}：

\begin{lstlisting}
error: config.yml: unknown key `layout.fnot-size` (see schools/default.yml)
\end{lstlisting}

所有可用的 key 與預設值都寫在 \code{schools/default.yml}，每一項都有註解說明單位，該份檔案就是完整的設定清單。

\section{主要檔案}

紙本論文有固定的裝訂順序：封面、審定書、目次、致謝、摘要、正文、參考文獻、附錄。\code{main.tex} 就是這份順序表。它不放任何論文內容，只寫明哪一個檔案排在哪個位置，以及從哪裡開始編章號。

換句話說，內容寫在 \code{contents/} 的各個檔案裡，\code{main.tex} 決定它們的先後；要調整裝訂順序、增刪某一頁，都是改這個檔案，不必去動內容。整份檔案很短，開頭到結尾就是下面這些：

\begin{lstlisting}[language=latex]
\documentclass{tw-thesis}

\begin{document}

\thesisfront

% front matter
\input{contents/front/acknowledgement}
\input{contents/front/abstract}

% notation
\input{contents/front/denotation}

% chapters
\input{contents/chapter01}
\input{contents/chapter02}

% references
\references

% appendices
\appendix
\input{contents/back/appendix01}

\end{document}
\end{lstlisting}

每一行的作用整理成\tblref{tbl:main}。

\begin{tlt}[notes={封面、審定書與三份目錄由 \texttt{\textbackslash thesisfront} 產生。}]{\texttt{main.tex} 各行的作用}{tbl:main}{L{0.5\linewidth}L{0.4\linewidth}}
內容 & 作用 \\
\midrule
\code{\documentclass{tw-thesis}} & 載入模板，讀入設定，以及圖片、表格、程式碼、假文等工具 \\
\code{\thesisfront} & 產生封面、審定書、目次、圖次、表次 \\
\code{\input{contents/front/...}} & 前置各頁，排版順序即此檔的先後順序 \\
\code{\input{contents/chapterNN}} & 正文各章，\code{\input} 的順序即為章號 \\
\code{\references} & 參考文獻，可指定文獻管理檔案位置 \\
\code{\appendix} & 此行後的章標題改用附錄編號 \code{A}、\code{B} \\
\code{\input{contents/back/...}} & 附錄 \\
\end{tlt}

幾件事需要特別注意：

\begin{enumerate}
\item 編號分界：致謝、摘要、符號列表用的是不編號的標題，所以不會變成第一章；後面用 \code{\chapter} 寫的章才會編號。
\item 章號看的是 \code{\input} 的排列順序，不是檔名，例如將 \code{chapter03} 那行移到 \code{chapter02} 前面，就會變成第二章。
\item 檔名雖然不影響結果，還是建議跟順序保持一致，免得日後自己找不到檔案。
\end{enumerate}

\section{新增與刪除章節}

模板把內容切成一章一個檔案，全部放在 \code{contents/}，檔名依章次編號，例如 \code{chapter01.tex}、\code{chapter02.tex}。改某一章不會動到其他章，編譯出錯時也可以快速判斷是哪個檔案的問題。

\code{main.tex} 裡的 \code{\input} 決定了各章的順序，新增或刪除章節時都需要自己維護這幾行。

新增一章：建立 \code{contents/chapter05.tex}（檔名依章次編號），檔案開頭放一個章標題與一個節標題，接著到 \code{main.tex} 的 \code{\references} 之前加入 \code{\input{contents/chapter05}}。

\begin{lstlisting}[language=latex]
\chapter{章標題}

\section{節標題}

內文。
\end{lstlisting}

刪除一章：把檔案刪掉，同時移除 \code{main.tex} 裡對應的那行 \code{\input}，否則 \LaTeX{} 找不到檔案會直接編譯失敗。後面各章的章號會自動遞補，正文裡的 \code{\ref} 引用也會跟著更新，都不用手改。

前置部分（\code{contents/front/}）與附錄（\code{contents/back/}）的做法相同：建立檔案，並在 \code{main.tex} 對應位置加一行 \code{\input}；刪除時也要一併移除那行。

\section{撰寫章節內容}

章節檔案建好後即可開始撰寫論文。建議先把章標題與幾個節標題列出來，再逐節填內文，不必一開始就寫完。

一章最基本的樣子如下：

\begin{lstlisting}[language=latex]
\chapter{緒論}

\section{研究背景}

第一段內文。

第二段內文，段落之間空一行。

\subsection{研究動機}

第三層標題底下的內文。
\end{lstlisting}

\code{\chapter}、\code{\section}、\code{\subsection} 依序是章、節、小節，第四層用 \code{\subsubsection}，不編號。每一層的字級、對齊方式與編號用詞都由學校設定檔決定，寫作時只要選對層級，不必自己調整外觀。章標題還會自動換頁，所以檔案開頭與結尾都不需要留空白。

段落之間空一行就會分段，首行縮排也是自動的。不要在每行末尾手動斷行，\LaTeX{} 會依版面寬度自行折行；原始碼裡的單一換行只會被當成一個空白，強制斷行反而會讓日後修改文字時出現奇怪的排版。

要讓某一章能被引用，在標題後面加上標籤，例如 \code{\chapter{研究方法}\label{cha:method}}。之後在任何地方寫 \code{\ref{cha:method}}，就會顯示成可點擊的「第三章」，章號變動時也會自動更新。圖、表與程式碼的標籤寫法相同，詳見第三章。

\section{暫時用假文看版面}

章節架構先定好、內文還沒寫時，可以先填假文，把版面、分頁與圖表位置排出來。\code{tw-thesis.sty} 包了兩個指令，預設各給五段：

\begin{lstlisting}[language=latex]
\zhlorem        % 中文五段
\enlorem        % 英文五段
\zhlorem[2]     % 顯示兩段
\enlorem[8]     % 顯示八段
\end{lstlisting}

中文假文來自 \href{https://ctan.org/pkg/zhlipsum}{zhlipsum} 套件，英文假文來自 \href{https://ctan.org/pkg/lipsum}{lipsum} 套件。兩個套件都隨 \TeX{} 發行版安裝。想換成別的文本或調整生成方式，改 \code{tw-thesis.sty} 檔尾的那兩行即可。

\chapter{版面設定}\label{cha:layout}

完成第一章的設定後，請先編譯一次，對照系所要求檢查字級、行距與頁面邊界，再調整需要修改的項目。本章會依序說明這些設定的位置與用法，建議一次修改一項，確認排版結果後再繼續。

下面的 \code{YAML} 範例都可以直接寫進 \code{config.yml} 覆蓋。若檔案裡已經有 \code{layout:}，請將新增的設定放在原有區塊下，保留範例中的縮排，不要再建立第二個同名區塊。每次修改設定後，都要重新執行 \code{python3 scripts/gen.py}（或直接用 \code{make}），\LaTeX{} 才讀得到新的值。

\section{字體與字級}

請先打開 \code{config.yml}，在 \code{layout:} 下設定本文的字型與字級。以下列出預設使用的字型；若目前顯示正常，可以先保留這些值，只調整系所要求的項目。

\begin{lstlisting}[language=yaml]
layout:
  language: "chinese" # chinese | english 決定標題語言
  font-latin: ["EB Garamond", "Tinos"] # 可填一個或多個，依序尋找
  font-cjk: ["TW-MOE-Std-Kai", "DFKai-SB"]
  font-size: 12 # pt
\end{lstlisting}

模板會依字型清單的順序尋找，用第一個找得到的字型：英文優先使用 EB Garamond，找不到才改用 Tinos（Tinos 是字寬與 Times New Roman 完全相同的開源字型），中文則優先使用\href{https://language.moe.gov.tw/material/info?m=9fe3fe82-8bbf-44c0-961d-873ea079e284}{教育部標準楷書}（授權條款為創用 CC 姓名標示─禁止改作，引用時須標示「中華民國教育部」）。這些字型檔都已經附在專案的 \code{fonts/} 裡，由 \code{tw-thesis.cls} 直接依檔案載入，不必另外安裝。兩套英文字型的授權都是 SIL OFL 1.1，可以自由散布。若系所明確要求 Times New Roman，可以把 \code{font-latin} 改成 \code{"Tinos"}：它的字寬與 Times New Roman 逐字相同，字形也相近，分頁位置不會跑掉；若一定要用原字型，因它有散布限制、模板無法附上，需先在自己的電腦安裝，再填入 \code{"Times New Roman"}。EB Garamond 的字寬與 Times New Roman 不同（平均約差 7\%），換過去後整份論文的分頁位置會改變。

若要換成系統或系所指定的其他字型，把字型名稱填進 \code{font-latin} 或 \code{font-cjk} 就好；不在 \code{fonts/} 裡的字型，需要已經安裝在系統上。這裡容易卡住的地方是：填的名稱要和字型檔內記載的家族名稱一致，而不是安裝畫面或字型網站上顯示的名稱，兩者有時不同。\tblref{tbl:font-name}列出幾種常見情境下查詢家族名稱的方式。

\begin{tlt}[notes={不確定時，一律以 \texttt{fc-list} 印出的名稱為準。}]{查詢字型家族名稱}{tbl:font-name}{L{0.28\linewidth}L{0.62\linewidth}}
情境 & 查詢方式 \\
\midrule
已經安裝在系統的字型 & 執行 \code{fc-list : family}，印出的名稱就是 \XeLaTeX{} 實際讀到的家族名，直接填入設定檔最準確。 \\
macOS 已安裝的字型 & 開啟字體簿（Font Book），點選字型後即可看到名稱。 \\
Windows 已安裝的字型 & 開啟「設定 → 個人化 → 字型」，點進字型頁面即可看到名稱。 \\
從 Google Fonts 等網站下載 & 頁面上顯示的名稱通常就是家族名，但下載後需先安裝到系統，再用 \code{fc-list} 確認一次。 \\
\end{tlt}

若要把本文改成 14 點，請在 \code{layout:} 下將 \code{font-size} 改為 \code{14}，不需要在數字後面加上單位。修改後可以找一頁中英混排的內文，檢查兩種文字的大小是否合適。章節標題有各自的字級設定，需到後面的 \code{heading:} 區塊另外調整。

\section{行高與段落}

本文的行距、段距、首行縮排與左右對齊都放在 \code{layout:} 下。可以先挑一頁有數個完整段落的內文，套用下面的設定，再比較行與行、段與段之間的距離，以及各段開頭的位置。

\begin{lstlisting}[language=yaml]
layout:
  line-height: 1.5      # 行高 1.5 倍
  par-spacing: ~        # ~ 跟隨行高；也可填 em 數值獨立調整
  first-line-indent: 2  # em，中文習慣首行縮排兩個字
  justify: true         # 左右對齊
\end{lstlisting}

調整段落時，可以先從 \code{line-height} 與 \code{first-line-indent} 開始。本文為 12 點、行高設為 \code{1.5} 時，一般文字行的基線間距為 18 點；首行縮排設為 \code{2}，則會留出兩個字寬的空間。若系所要求段落不縮排，將 \code{first-line-indent} 改為 \code{0} 即可。

\code{par-spacing} 控制兩個段落之間的間距。預設的 \code{~} 會跟隨 \code{line-height}，讓換段前後的基線距離與一般行距相同；例如 \code{line-height: 1.5} 會換算成 \code{0.5em} 的段距，加上 1em 的文字行框後，兩段之間仍維持 1.5 倍行高，不會額外多出一行空白。

若學校或系所另有規定，可以將 \code{par-spacing} 改成數字獨立覆寫，單位是 \code{em}，不需在 \code{YAML} 中加上單位；例如 \code{par-spacing: 1.2} 會加大段落間的留白。這個設定不會改變 \code{line-height}，之後調整行距時也不會跟著變動。章節原始碼中以空白行分隔的文字會形成不同段落，可以用相鄰的兩段文字觀察調整結果。

\code{config.yml} 裡還有一個 \code{line-box}，它是 Typst 版用來釘住行盒上下緣的設定；\LaTeX{} 直接以基線間距排版，用不到這一項，填了也不會有作用，保留只是為了讓兩個版本能共用同一份設定檔。

\section{邊界與紙張}

紙張大小由 \code{paper} 決定，四邊留白則在 \code{margin:} 下分別設定。以下使用 A4 紙張；若已選擇學校格式，請先查看套用結果，只加入需要覆蓋的邊界數值即可。

\begin{lstlisting}[language=yaml]
layout:
  paper: "a4"
  margin: # cm
    top: 3
    bottom: 2
    left: 3
    right: 3
\end{lstlisting}

請在 \code{margin:} 下填入四個方向的邊界，單位為公分，直接寫數字即可。例如系所要求左側留 3.5 公分，就將 \code{left} 設為 \code{3.5}，不要寫成 \code{3.5cm}。若只修改左側，其餘三邊仍會沿用學校設定或預設值，不必全部重寫。

修改後，請同時查看正文與封面，確認標題、圖表和封面資料是否仍放得下。封面的邊界由 \code{cover.margin} 另外設定，沒有填的邊會沿用正文邊界；封面各區塊另以固定間距與彈性間距安排位置。若論文題目較長，可以再調整 \code{cover:} 下的字級或間距，避免文字擠在一起。

\section{標題與編號}

章節標題的字級、對齊方式與前後間距，都集中在 \code{layout.heading} 區。下列各組清單依照標題層級排列，可以先對照章、節與小節的外觀，再調整對應位置的數值。

\begin{lstlisting}[language=yaml]
layout:
  heading:
    size: [18, 16, 14, 12]                     # pt，第 1 至 4 層
    weight: "regular"                          # regular | bold
    align: ["center", "center", "left", "left"]
    line-height: 1.2                           # 標題內部的行距
    gap-above: [1, 0.8, 0.8, 0.6]              # em
    gap-below: [1, 0.5, 0.4, 0.3]              # em
    number-gap: [1, 1, 0, 0]                   # em，編號與標題之間
\end{lstlisting}

上面的清單依序對應第一至第四層標題。例如要調整章標題的字級，就修改 \code{size} 的第一個數值；要調整節標題，則修改第二個數值。建議先保留四個位置，只更換需要的值，方便對照各層的設定。

若要修改編號前後的文字，請在 \code{config.yml} 另外加入 \code{strings:}，與 \code{layout:} 放在同一層。以下範例會讓章、節分別使用「第一章」與「第一節」的寫法：

\begin{lstlisting}[language=yaml]
strings:
  chapter: ["第", "章"]   # 第一章
  section: ["第", "節"]   # 第一節
  subsection: "、"        # 一、
  appendix: ["附錄 ", ""] # 附錄 A
\end{lstlisting}

回到章節檔案後，用 \LaTeX{} 的分節指令指定層級。\code{\chapter} 是章標題，\code{\section} 是節標題，依此類推。標題文字寫在大括號內，並依照內容的層次安排：

\begin{lstlisting}[language=latex]
\chapter{研究方法}       % 第一層，章號依出現順序自動產生
\section{實驗設計}       % 第二層，印成第一節
\subsection{資料來源}    % 第三層，印成一、
\subsubsection{前處理}   % 第四層，不編號
\end{lstlisting}

編譯時，第一層標題會自動另起一頁，章號也會依出現順序產生。因此標題只需寫「研究方法」，不必手動加上「第三章」或插入換頁指令。新增或調整章節順序後，重新編譯就會更新編號；因為目次與交叉引用要靠編號檔記錄，\LaTeX{} 需要連續編譯兩次才會完全正確，\code{latexmk} 會自動處理。

\section{圖表與程式碼的標題位置}\label{sec:caption}

圖片、表格與程式碼的標題位置可以分別設定。請在 \code{config.yml} 的 \code{layout:} 下找到 \code{caption-position:}，將要調整的項目設為 \code{"top"} 或 \code{"bottom"}，分別代表放在內容上方或下方。以下是目前的預設值：

\begin{lstlisting}[language=yaml]
layout:
  caption-position:
    image: "bottom"    # 圖片標題放在下方
    table: "bottom"    # 表格標題放在下方
    code: "bottom"     # 程式碼標題放在下方
\end{lstlisting}

例如要將所有表格的標題移到上方，只需把 \code{table} 改為 \code{"top"}，圖片與程式碼不會跟著改變。圖片的設定也會套用到並排與網格圖片的總標題。修改後重新編譯，檢查標題與內容的位置即可，不需要逐一修改各章的圖表。

\section{目次與頁碼}

目次的收錄層級、縮排與頁碼格式都放在 \code{outline:} 下，與 \code{layout:} 位於同一層。以下列出常用設定，初次使用時可以先保留預設值，等章節內容齊全後再調整。

\begin{lstlisting}[language=yaml]
outline:
  depth: 3            # 目次收到第幾層
  indent: 1           # em，每層縮排一個字
  title-size: 16      # pt，目次、圖次、表次的字級
  page-numbering: "i" # 前置部分：i, ii, iii
  body-numbering: "1" # 正文開始：1, 2, 3
  entry-gap: 1        # em，章與章之間的間距
  entry-bold: true    # 章標題在目次裡加粗
\end{lstlisting}

目次會依照章節標題自動產生，圖次與表次則整理文件中的圖片和表格。若只想在目次中列出章與節，將 \code{depth} 改為 \code{2} 即可，正文的標題仍會保留原有層級。編譯後可以點擊目次條目，確認是否跳到對應位置。

要修改這三頁的標題，請在 \code{titles:} 下設定文字。字串中的 \code{~} 會轉成字距，例如 \code{"目~次"} 會拉開兩個字的間隔；若要直接顯示「目錄」，填入 \code{"目錄"} 即可：

\begin{lstlisting}[language=yaml]
titles:
  toc: "目~次"   # 有些學校用「目錄」
  lof: "圖~次"
  lot: "表~次"
\end{lstlisting}

頁碼的位置與顯示文字放在 \code{page-number:} 下，與 \code{outline:} 位於同一層。以下都是預設值，不修改也能直接使用：

\begin{lstlisting}[language=yaml]
page-number:
  position: "bottom" # top：頁首；bottom：頁尾
  align: "center"    # left：靠左；center：置中；right：靠右
  format: "{page}"   # {page} 會替換成該頁的頁碼
\end{lstlisting}

各項設定可依下表調整。將 \code{position} 設為 \code{"top"}、\code{align} 設為 \code{"right"}，頁碼就會移到頁首靠右。

\begin{tlt}{頁碼設定參數}{tbl:page-number}{L{0.15\linewidth}L{0.75\linewidth}}
參數 & 用途與範例 \\
\midrule
\code{position} & 位置：\code{"top"} 頁首、\code{"bottom"} 頁尾（預設）。 \\
\code{align} & 對齊：\code{"left"}、\code{"center"}（預設）、\code{"right"}。 \\
\code{format} & 顯示文字：\code{"第 {page} 頁"}；設為 \code{""} 可隱藏頁碼。 \\
\end{tlt}

\code{format} 中的 \code{{page}} 會替換成頁碼，例如 \code{"第 {page} 頁"} 顯示為「第 1 頁」；也可以寫成 \code{"論文 {page}"}。其他文字會直接顯示，特殊符號（如 \code{%}、\code{&}）由產生設定的程式自動處理，不需要加反斜線。

\code{outline.page-numbering} 與 \code{outline.body-numbering} 分別控制前置部分與正文的編號方式，預設為小寫羅馬數字與阿拉伯數字。\code{page-number.format} 只調整頁首或頁尾的文字，目次仍保留原本的頁碼格式。封面與審定書不顯示頁碼。完成章節調整後，請檢查正文第一頁與目次所列的頁碼，確認兩者對應正確。

\section{浮水印與定稿}

提交前若需要加上浮水印，請先檢查學校設定是否已指定圖片。需要自行指定時，可以在 \code{config.yml} 加入下列區塊，將 \code{image} 改成要使用的圖檔路徑，再調整大小與深淺。

\begin{lstlisting}[language=yaml]
watermark:
  image: "/assets/watermarks/nccu.pdf"
  width: 100    # %，佔頁寬的比例，置中
  opacity: 100  # %，100 是原圖深淺，50 會刷淡一半
\end{lstlisting}

設定好圖片後，執行 \code{make release}，再打開 PDF 的正文頁面查看浮水印。一般的 \code{make} 不會啟用浮水印，封面與前置頁面也不會加上。若 \code{image} 的值是 \code{~}，表示沒有指定圖片，即使執行 \code{make release} 也不會顯示。

若要換成自己的浮水印，請先將圖片放進 \code{assets/watermarks/}，再修改 \code{image} 的路徑。可以用 \code{width} 調整圖片占頁寬的比例，或降低 \code{opacity}，讓浮水印淡一些。是否需要自行加上浮水印，請依提交時的學校要求處理。

沒有安裝 \code{make} 時，也可以在專案根目錄執行以下指令，輸出的檔案同樣是 \code{main.pdf}：

\begin{lstlisting}[language=bash]
python3 scripts/gen.py
latexmk -g -xelatex -usepretex='\def\twwatermark{1}' main.tex
\end{lstlisting}


% references
\references

% appendices
\appendix
\chapter{編譯問題排查}\label{appendix-a}

使用模板時，若遇到編譯失敗、圖片找不到或 PDF 沒有更新，可以先依本附錄的順序檢查。建議保留終端機中的錯誤訊息，從第一個錯誤指出的檔案與位置開始處理，修正後再重新編譯。

這份檔案也示範附錄的排版方式。換成自己的論文時，可以保留章標題的寫法，將標題與內文替換成問卷、補充數據或其他資料，再到 \code{main.tex} 確認附錄的裝訂順序。

\section{指令執行失敗}

若終端機顯示找不到 \code{xelatex} 或 \code{latexmk}，請先確認已完成安裝，再重新開啟終端機，執行

\begin{lstlisting}[language=bash]
xelatex --version
\end{lstlisting}

確認能夠顯示版本資訊後，再嘗試編譯。若找不到的是 \code{make}，則可以先使用以下指令，直接產生 PDF：

\begin{lstlisting}[language=bash]
texlua scripts/gen.lua
latexmk -xelatex main.tex
\end{lstlisting}

若錯誤訊息指出找不到 \code{main.tex} 或 \code{.tw-thesis-config.tex}，請檢查終端機目前所在的資料夾。編譯指令應在專案根目錄執行，也就是能看到 \code{main.tex}、\code{config.yml} 與 \code{contents/} 的位置；\code{.tw-thesis-config.tex} 則是 \code{scripts/gen.lua} 產生的，先執行過一次就會出現。若目前位於章節資料夾，請先切換回根目錄再試一次。

\section{修改設定後無法編譯}

修改設定後出現錯誤時，請先看訊息指出哪個檔案、哪一個欄位。若包含 \code{unknown key}，可以對照 \code{schools/default.yml} 檢查欄位拼字與所在層級；若是 \code{YAML} 解析錯誤，則先確認縮排使用空白，且冒號後保留空白，再檢查引號是否成對。

加入新設定時，請先找出檔案中是否已有相同區塊，避免在同一層重複建立 \code{layout:}。例如要同時修改字級與左側邊界，應把兩項設定放在同一個 \code{layout:} 下，並讓 \code{left} 再往內縮排一層：

\begin{lstlisting}[language=yaml]
layout:
  font-size: 12
  margin:
    left: 3.5
\end{lstlisting}

若一次修改了多個設定，暫時無法判斷錯誤來源，可以先將剛加入的內容複製到別處保存，再移除這些設定，確認原本的文件能正常編譯。之後逐項加回，每次存檔後檢查結果，較容易找到需要修正的位置。

\section{圖片或引用找不到}

圖片無法載入時，請先確認檔案確實放在 \code{figures/} 的指定位置，再比對路徑中的檔名、副檔名與大小寫。路徑是從專案根目錄算起的，不是從目前章節所在的資料夾開始。若曾經移動或重新命名圖片，也要一併修改章節中的路徑。

圖表引用顯示成 \code{??} 時，請檢查圖表是否已加上 \code{\label}，以及名稱是否與內文的 \code{\ref} 完全相同；也別忘了 \LaTeX{} 需要編譯兩次才會解出交叉引用，\code{latexmk} 會自動處理。若是文獻引用出錯，則到 \code{.bib} 檔案確認識別名稱，再檢查 \code{config.yml} 的 \code{bibliography.path} 是否包含該文獻檔。修正後重新編譯，確認引用已顯示為編號或文獻資料。

\section{PDF 沒有更新}

修改內容後，請先確認終端機是否已完成這次編譯，以及是否出現錯誤。編譯失敗時，資料夾中可能仍保留上一次的 \code{main.pdf}，因此能開啟 PDF 並不代表新內容已經寫入。請先修正錯誤，再重新執行編譯指令。

若使用 \code{make watch}，請確認監看程序仍在執行，並且修改的檔案已經存檔。編譯成功後畫面仍未更新時，可以關閉再重新開啟 PDF，同時確認檢視器開啟的是目前專案中的 \code{main.pdf}，避免看到先前複製到其他位置的版本。若編譯結果卡在舊的狀態，執行 \code{make clean} 清掉暫存檔後再編譯一次，通常就能解決。


\end{document}
\end{lstlisting}

每一行的作用整理成\tblref{tbl:main}。

\begin{tlt}[notes={封面、審定書與三份目錄由 \texttt{\textbackslash thesisfront} 產生。}]{\texttt{main.tex} 各行的作用}{tbl:main}{L{0.5\linewidth}L{0.4\linewidth}}
內容 & 作用 \\
\midrule
\code{\documentclass{tw-thesis}} & 載入模板，讀入設定，以及圖片、表格、程式碼、假文等工具 \\
\code{\thesisfront} & 產生封面、審定書、目次、圖次、表次 \\
\code{\input{contents/front/...}} & 前置各頁，排版順序即此檔的先後順序 \\
\code{\input{contents/chapterNN}} & 正文各章，\code{\input} 的順序即為章號 \\
\code{\references} & 參考文獻，可指定文獻管理檔案位置 \\
\code{\appendix} & 此行後的章標題改用附錄編號 \code{A}、\code{B} \\
\code{\input{contents/back/...}} & 附錄 \\
\end{tlt}

幾件事需要特別注意：

\begin{enumerate}
\item 編號分界：致謝、摘要、符號列表用的是不編號的標題，所以不會變成第一章；後面用 \code{\chapter} 寫的章才會編號。
\item 章號看的是 \code{\input} 的排列順序，不是檔名，例如將 \code{chapter03} 那行移到 \code{chapter02} 前面，就會變成第二章。
\item 檔名雖然不影響結果，還是建議跟順序保持一致，免得日後自己找不到檔案。
\end{enumerate}

\section{新增與刪除章節}

模板把內容切成一章一個檔案，全部放在 \code{contents/}，檔名依章次編號，例如 \code{chapter01.tex}、\code{chapter02.tex}。改某一章不會動到其他章，編譯出錯時也可以快速判斷是哪個檔案的問題。

\code{main.tex} 裡的 \code{\input} 決定了各章的順序，新增或刪除章節時都需要自己維護這幾行。

新增一章：建立 \code{contents/chapter05.tex}（檔名依章次編號），檔案開頭放一個章標題與一個節標題，接著到 \code{main.tex} 的 \code{\references} 之前加入 \code{\input{contents/chapter05}}。

\begin{lstlisting}[language=latex]
\chapter{章標題}

\section{節標題}

內文。
\end{lstlisting}

刪除一章：把檔案刪掉，同時移除 \code{main.tex} 裡對應的那行 \code{\input}，否則 \LaTeX{} 找不到檔案會直接編譯失敗。後面各章的章號會自動遞補，正文裡的 \code{\ref} 引用也會跟著更新，都不用手改。

前置部分（\code{contents/front/}）與附錄（\code{contents/back/}）的做法相同：建立檔案，並在 \code{main.tex} 對應位置加一行 \code{\input}；刪除時也要一併移除那行。

\section{撰寫章節內容}

章節檔案建好後即可開始撰寫論文。建議先把章標題與幾個節標題列出來，再逐節填內文，不必一開始就寫完。

一章最基本的樣子如下：

\begin{lstlisting}[language=latex]
\chapter{緒論}

\section{研究背景}

第一段內文。

第二段內文，段落之間空一行。

\subsection{研究動機}

第三層標題底下的內文。
\end{lstlisting}

\code{\chapter}、\code{\section}、\code{\subsection} 依序是章、節、小節，第四層用 \code{\subsubsection}，不編號。每一層的字級、對齊方式與編號用詞都由學校設定檔決定，寫作時只要選對層級，不必自己調整外觀。章標題還會自動換頁，所以檔案開頭與結尾都不需要留空白。

段落之間空一行就會分段，首行縮排也是自動的。不要在每行末尾手動斷行，\LaTeX{} 會依版面寬度自行折行；原始碼裡的單一換行只會被當成一個空白，強制斷行反而會讓日後修改文字時出現奇怪的排版。

要讓某一章能被引用，在標題後面加上標籤，例如 \code{\chapter{研究方法}\label{cha:method}}。之後在任何地方寫 \code{\ref{cha:method}}，就會顯示成可點擊的「第三章」，章號變動時也會自動更新。圖、表與程式碼的標籤寫法相同，詳見第三章。

\section{暫時用假文看版面}

章節架構先定好、內文還沒寫時，可以先填假文，把版面、分頁與圖表位置排出來。\code{tw-thesis.sty} 包了兩個指令，預設各給五段：

\begin{lstlisting}[language=latex]
\zhlorem        % 中文五段
\enlorem        % 英文五段
\zhlorem[2]     % 顯示兩段
\enlorem[8]     % 顯示八段
\end{lstlisting}

中文假文來自 \href{https://ctan.org/pkg/zhlipsum}{zhlipsum} 套件，英文假文來自 \href{https://ctan.org/pkg/lipsum}{lipsum} 套件。兩個套件都隨 \TeX{} 發行版安裝。想換成別的文本或調整生成方式，改 \code{tw-thesis.sty} 檔尾的那兩行即可。

\chapter{版面設定}\label{cha:layout}

完成第一章的設定後，請先編譯一次，對照系所要求檢查字級、行距與頁面邊界，再調整需要修改的項目。本章會依序說明這些設定的位置與用法，建議一次修改一項，確認排版結果後再繼續。

下面的 \code{YAML} 範例都可以直接寫進 \code{config.yml} 覆蓋。若檔案裡已經有 \code{layout:}，請將新增的設定放在原有區塊下，保留範例中的縮排，不要再建立第二個同名區塊。每次修改設定後，都要重新執行 \code{python3 scripts/gen.py}（或直接用 \code{make}），\LaTeX{} 才讀得到新的值。

\section{字體與字級}

請先打開 \code{config.yml}，在 \code{layout:} 下設定本文的字型與字級。以下列出預設使用的字型；若目前顯示正常，可以先保留這些值，只調整系所要求的項目。

\begin{lstlisting}[language=yaml]
layout:
  language: "chinese" # chinese | english 決定標題語言
  font-latin: ["EB Garamond", "Tinos"] # 可填一個或多個，依序尋找
  font-cjk: ["TW-MOE-Std-Kai", "DFKai-SB"]
  font-size: 12 # pt
\end{lstlisting}

模板會依字型清單的順序尋找，用第一個找得到的字型：英文優先使用 EB Garamond，找不到才改用 Tinos（Tinos 是字寬與 Times New Roman 完全相同的開源字型），中文則優先使用\href{https://language.moe.gov.tw/material/info?m=9fe3fe82-8bbf-44c0-961d-873ea079e284}{教育部標準楷書}（授權條款為創用 CC 姓名標示─禁止改作，引用時須標示「中華民國教育部」）。這些字型檔都已經附在專案的 \code{fonts/} 裡，由 \code{tw-thesis.cls} 直接依檔案載入，不必另外安裝。兩套英文字型的授權都是 SIL OFL 1.1，可以自由散布。若系所明確要求 Times New Roman，可以把 \code{font-latin} 改成 \code{"Tinos"}：它的字寬與 Times New Roman 逐字相同，字形也相近，分頁位置不會跑掉；若一定要用原字型，因它有散布限制、模板無法附上，需先在自己的電腦安裝，再填入 \code{"Times New Roman"}。EB Garamond 的字寬與 Times New Roman 不同（平均約差 7\%），換過去後整份論文的分頁位置會改變。

若要換成系統或系所指定的其他字型，把字型名稱填進 \code{font-latin} 或 \code{font-cjk} 就好；不在 \code{fonts/} 裡的字型，需要已經安裝在系統上。這裡容易卡住的地方是：填的名稱要和字型檔內記載的家族名稱一致，而不是安裝畫面或字型網站上顯示的名稱，兩者有時不同。\tblref{tbl:font-name}列出幾種常見情境下查詢家族名稱的方式。

\begin{tlt}[notes={不確定時，一律以 \texttt{fc-list} 印出的名稱為準。}]{查詢字型家族名稱}{tbl:font-name}{L{0.28\linewidth}L{0.62\linewidth}}
情境 & 查詢方式 \\
\midrule
已經安裝在系統的字型 & 執行 \code{fc-list : family}，印出的名稱就是 \XeLaTeX{} 實際讀到的家族名，直接填入設定檔最準確。 \\
macOS 已安裝的字型 & 開啟字體簿（Font Book），點選字型後即可看到名稱。 \\
Windows 已安裝的字型 & 開啟「設定 → 個人化 → 字型」，點進字型頁面即可看到名稱。 \\
從 Google Fonts 等網站下載 & 頁面上顯示的名稱通常就是家族名，但下載後需先安裝到系統，再用 \code{fc-list} 確認一次。 \\
\end{tlt}

若要把本文改成 14 點，請在 \code{layout:} 下將 \code{font-size} 改為 \code{14}，不需要在數字後面加上單位。修改後可以找一頁中英混排的內文，檢查兩種文字的大小是否合適。章節標題有各自的字級設定，需到後面的 \code{heading:} 區塊另外調整。

\section{行高與段落}

本文的行距、段距、首行縮排與左右對齊都放在 \code{layout:} 下。可以先挑一頁有數個完整段落的內文，套用下面的設定，再比較行與行、段與段之間的距離，以及各段開頭的位置。

\begin{lstlisting}[language=yaml]
layout:
  line-height: 1.5      # 行高 1.5 倍
  par-spacing: ~        # ~ 跟隨行高；也可填 em 數值獨立調整
  first-line-indent: 2  # em，中文習慣首行縮排兩個字
  justify: true         # 左右對齊
\end{lstlisting}

調整段落時，可以先從 \code{line-height} 與 \code{first-line-indent} 開始。本文為 12 點、行高設為 \code{1.5} 時，一般文字行的基線間距為 18 點；首行縮排設為 \code{2}，則會留出兩個字寬的空間。若系所要求段落不縮排，將 \code{first-line-indent} 改為 \code{0} 即可。

\code{par-spacing} 控制兩個段落之間的間距。預設的 \code{~} 會跟隨 \code{line-height}，讓換段前後的基線距離與一般行距相同；例如 \code{line-height: 1.5} 會換算成 \code{0.5em} 的段距，加上 1em 的文字行框後，兩段之間仍維持 1.5 倍行高，不會額外多出一行空白。

若學校或系所另有規定，可以將 \code{par-spacing} 改成數字獨立覆寫，單位是 \code{em}，不需在 \code{YAML} 中加上單位；例如 \code{par-spacing: 1.2} 會加大段落間的留白。這個設定不會改變 \code{line-height}，之後調整行距時也不會跟著變動。章節原始碼中以空白行分隔的文字會形成不同段落，可以用相鄰的兩段文字觀察調整結果。

\code{config.yml} 裡還有一個 \code{line-box}，它是 Typst 版用來釘住行盒上下緣的設定；\LaTeX{} 直接以基線間距排版，用不到這一項，填了也不會有作用，保留只是為了讓兩個版本能共用同一份設定檔。

\section{邊界與紙張}

紙張大小由 \code{paper} 決定，四邊留白則在 \code{margin:} 下分別設定。以下使用 A4 紙張；若已選擇學校格式，請先查看套用結果，只加入需要覆蓋的邊界數值即可。

\begin{lstlisting}[language=yaml]
layout:
  paper: "a4"
  margin: # cm
    top: 3
    bottom: 2
    left: 3
    right: 3
\end{lstlisting}

請在 \code{margin:} 下填入四個方向的邊界，單位為公分，直接寫數字即可。例如系所要求左側留 3.5 公分，就將 \code{left} 設為 \code{3.5}，不要寫成 \code{3.5cm}。若只修改左側，其餘三邊仍會沿用學校設定或預設值，不必全部重寫。

修改後，請同時查看正文與封面，確認標題、圖表和封面資料是否仍放得下。封面的邊界由 \code{cover.margin} 另外設定，沒有填的邊會沿用正文邊界；封面各區塊另以固定間距與彈性間距安排位置。若論文題目較長，可以再調整 \code{cover:} 下的字級或間距，避免文字擠在一起。

\section{標題與編號}

章節標題的字級、對齊方式與前後間距，都集中在 \code{layout.heading} 區。下列各組清單依照標題層級排列，可以先對照章、節與小節的外觀，再調整對應位置的數值。

\begin{lstlisting}[language=yaml]
layout:
  heading:
    size: [18, 16, 14, 12]                     # pt，第 1 至 4 層
    weight: "regular"                          # regular | bold
    align: ["center", "center", "left", "left"]
    line-height: 1.2                           # 標題內部的行距
    gap-above: [1, 0.8, 0.8, 0.6]              # em
    gap-below: [1, 0.5, 0.4, 0.3]              # em
    number-gap: [1, 1, 0, 0]                   # em，編號與標題之間
\end{lstlisting}

上面的清單依序對應第一至第四層標題。例如要調整章標題的字級，就修改 \code{size} 的第一個數值；要調整節標題，則修改第二個數值。建議先保留四個位置，只更換需要的值，方便對照各層的設定。

若要修改編號前後的文字，請在 \code{config.yml} 另外加入 \code{strings:}，與 \code{layout:} 放在同一層。以下範例會讓章、節分別使用「第一章」與「第一節」的寫法：

\begin{lstlisting}[language=yaml]
strings:
  chapter: ["第", "章"]   # 第一章
  section: ["第", "節"]   # 第一節
  subsection: "、"        # 一、
  appendix: ["附錄 ", ""] # 附錄 A
\end{lstlisting}

回到章節檔案後，用 \LaTeX{} 的分節指令指定層級。\code{\chapter} 是章標題，\code{\section} 是節標題，依此類推。標題文字寫在大括號內，並依照內容的層次安排：

\begin{lstlisting}[language=latex]
\chapter{研究方法}       % 第一層，章號依出現順序自動產生
\section{實驗設計}       % 第二層，印成第一節
\subsection{資料來源}    % 第三層，印成一、
\subsubsection{前處理}   % 第四層，不編號
\end{lstlisting}

編譯時，第一層標題會自動另起一頁，章號也會依出現順序產生。因此標題只需寫「研究方法」，不必手動加上「第三章」或插入換頁指令。新增或調整章節順序後，重新編譯就會更新編號；因為目次與交叉引用要靠編號檔記錄，\LaTeX{} 需要連續編譯兩次才會完全正確，\code{latexmk} 會自動處理。

\section{圖表與程式碼的標題位置}\label{sec:caption}

圖片、表格與程式碼的標題位置可以分別設定。請在 \code{config.yml} 的 \code{layout:} 下找到 \code{caption-position:}，將要調整的項目設為 \code{"top"} 或 \code{"bottom"}，分別代表放在內容上方或下方。以下是目前的預設值：

\begin{lstlisting}[language=yaml]
layout:
  caption-position:
    image: "bottom"    # 圖片標題放在下方
    table: "bottom"    # 表格標題放在下方
    code: "bottom"     # 程式碼標題放在下方
\end{lstlisting}

例如要將所有表格的標題移到上方，只需把 \code{table} 改為 \code{"top"}，圖片與程式碼不會跟著改變。圖片的設定也會套用到並排與網格圖片的總標題。修改後重新編譯，檢查標題與內容的位置即可，不需要逐一修改各章的圖表。

\section{目次與頁碼}

目次的收錄層級、縮排與頁碼格式都放在 \code{outline:} 下，與 \code{layout:} 位於同一層。以下列出常用設定，初次使用時可以先保留預設值，等章節內容齊全後再調整。

\begin{lstlisting}[language=yaml]
outline:
  depth: 3            # 目次收到第幾層
  indent: 1           # em，每層縮排一個字
  title-size: 16      # pt，目次、圖次、表次的字級
  page-numbering: "i" # 前置部分：i, ii, iii
  body-numbering: "1" # 正文開始：1, 2, 3
  entry-gap: 1        # em，章與章之間的間距
  entry-bold: true    # 章標題在目次裡加粗
\end{lstlisting}

目次會依照章節標題自動產生，圖次與表次則整理文件中的圖片和表格。若只想在目次中列出章與節，將 \code{depth} 改為 \code{2} 即可，正文的標題仍會保留原有層級。編譯後可以點擊目次條目，確認是否跳到對應位置。

要修改這三頁的標題，請在 \code{titles:} 下設定文字。字串中的 \code{~} 會轉成字距，例如 \code{"目~次"} 會拉開兩個字的間隔；若要直接顯示「目錄」，填入 \code{"目錄"} 即可：

\begin{lstlisting}[language=yaml]
titles:
  toc: "目~次"   # 有些學校用「目錄」
  lof: "圖~次"
  lot: "表~次"
\end{lstlisting}

頁碼的位置與顯示文字放在 \code{page-number:} 下，與 \code{outline:} 位於同一層。以下都是預設值，不修改也能直接使用：

\begin{lstlisting}[language=yaml]
page-number:
  position: "bottom" # top：頁首；bottom：頁尾
  align: "center"    # left：靠左；center：置中；right：靠右
  format: "{page}"   # {page} 會替換成該頁的頁碼
\end{lstlisting}

各項設定可依下表調整。將 \code{position} 設為 \code{"top"}、\code{align} 設為 \code{"right"}，頁碼就會移到頁首靠右。

\begin{tlt}{頁碼設定參數}{tbl:page-number}{L{0.15\linewidth}L{0.75\linewidth}}
參數 & 用途與範例 \\
\midrule
\code{position} & 位置：\code{"top"} 頁首、\code{"bottom"} 頁尾（預設）。 \\
\code{align} & 對齊：\code{"left"}、\code{"center"}（預設）、\code{"right"}。 \\
\code{format} & 顯示文字：\code{"第 {page} 頁"}；設為 \code{""} 可隱藏頁碼。 \\
\end{tlt}

\code{format} 中的 \code{{page}} 會替換成頁碼，例如 \code{"第 {page} 頁"} 顯示為「第 1 頁」；也可以寫成 \code{"論文 {page}"}。其他文字會直接顯示，特殊符號（如 \code{%}、\code{&}）由產生設定的程式自動處理，不需要加反斜線。

\code{outline.page-numbering} 與 \code{outline.body-numbering} 分別控制前置部分與正文的編號方式，預設為小寫羅馬數字與阿拉伯數字。\code{page-number.format} 只調整頁首或頁尾的文字，目次仍保留原本的頁碼格式。封面與審定書不顯示頁碼。完成章節調整後，請檢查正文第一頁與目次所列的頁碼，確認兩者對應正確。

\section{浮水印與定稿}

提交前若需要加上浮水印，請先檢查學校設定是否已指定圖片。需要自行指定時，可以在 \code{config.yml} 加入下列區塊，將 \code{image} 改成要使用的圖檔路徑，再調整大小與深淺。

\begin{lstlisting}[language=yaml]
watermark:
  image: "/assets/watermarks/nccu.pdf"
  width: 100    # %，佔頁寬的比例，置中
  opacity: 100  # %，100 是原圖深淺，50 會刷淡一半
\end{lstlisting}

設定好圖片後，執行 \code{make release}，再打開 PDF 的正文頁面查看浮水印。一般的 \code{make} 不會啟用浮水印，封面與前置頁面也不會加上。若 \code{image} 的值是 \code{~}，表示沒有指定圖片，即使執行 \code{make release} 也不會顯示。

若要換成自己的浮水印，請先將圖片放進 \code{assets/watermarks/}，再修改 \code{image} 的路徑。可以用 \code{width} 調整圖片占頁寬的比例，或降低 \code{opacity}，讓浮水印淡一些。是否需要自行加上浮水印，請依提交時的學校要求處理。

沒有安裝 \code{make} 時，也可以在專案根目錄執行以下指令，輸出的檔案同樣是 \code{main.pdf}：

\begin{lstlisting}[language=bash]
python3 scripts/gen.py
latexmk -g -xelatex -usepretex='\def\twwatermark{1}' main.tex
\end{lstlisting}


% references
\references

% appendices
\appendix
\chapter{編譯問題排查}\label{appendix-a}

使用模板時，若遇到編譯失敗、圖片找不到或 PDF 沒有更新，可以先依本附錄的順序檢查。建議保留終端機中的錯誤訊息，從第一個錯誤指出的檔案與位置開始處理，修正後再重新編譯。

這份檔案也示範附錄的排版方式。換成自己的論文時，可以保留章標題的寫法，將標題與內文替換成問卷、補充數據或其他資料，再到 \code{main.tex} 確認附錄的裝訂順序。

\section{指令執行失敗}

若終端機顯示找不到 \code{xelatex} 或 \code{latexmk}，請先確認已完成安裝，再重新開啟終端機，執行

\begin{lstlisting}[language=bash]
xelatex --version
\end{lstlisting}

確認能夠顯示版本資訊後，再嘗試編譯。若找不到的是 \code{make}，則可以先使用以下指令，直接產生 PDF：

\begin{lstlisting}[language=bash]
texlua scripts/gen.lua
latexmk -xelatex main.tex
\end{lstlisting}

若錯誤訊息指出找不到 \code{main.tex} 或 \code{.tw-thesis-config.tex}，請檢查終端機目前所在的資料夾。編譯指令應在專案根目錄執行，也就是能看到 \code{main.tex}、\code{config.yml} 與 \code{contents/} 的位置；\code{.tw-thesis-config.tex} 則是 \code{scripts/gen.lua} 產生的，先執行過一次就會出現。若目前位於章節資料夾，請先切換回根目錄再試一次。

\section{修改設定後無法編譯}

修改設定後出現錯誤時，請先看訊息指出哪個檔案、哪一個欄位。若包含 \code{unknown key}，可以對照 \code{schools/default.yml} 檢查欄位拼字與所在層級；若是 \code{YAML} 解析錯誤，則先確認縮排使用空白，且冒號後保留空白，再檢查引號是否成對。

加入新設定時，請先找出檔案中是否已有相同區塊，避免在同一層重複建立 \code{layout:}。例如要同時修改字級與左側邊界，應把兩項設定放在同一個 \code{layout:} 下，並讓 \code{left} 再往內縮排一層：

\begin{lstlisting}[language=yaml]
layout:
  font-size: 12
  margin:
    left: 3.5
\end{lstlisting}

若一次修改了多個設定，暫時無法判斷錯誤來源，可以先將剛加入的內容複製到別處保存，再移除這些設定，確認原本的文件能正常編譯。之後逐項加回，每次存檔後檢查結果，較容易找到需要修正的位置。

\section{圖片或引用找不到}

圖片無法載入時，請先確認檔案確實放在 \code{figures/} 的指定位置，再比對路徑中的檔名、副檔名與大小寫。路徑是從專案根目錄算起的，不是從目前章節所在的資料夾開始。若曾經移動或重新命名圖片，也要一併修改章節中的路徑。

圖表引用顯示成 \code{??} 時，請檢查圖表是否已加上 \code{\label}，以及名稱是否與內文的 \code{\ref} 完全相同；也別忘了 \LaTeX{} 需要編譯兩次才會解出交叉引用，\code{latexmk} 會自動處理。若是文獻引用出錯，則到 \code{.bib} 檔案確認識別名稱，再檢查 \code{config.yml} 的 \code{bibliography.path} 是否包含該文獻檔。修正後重新編譯，確認引用已顯示為編號或文獻資料。

\section{PDF 沒有更新}

修改內容後，請先確認終端機是否已完成這次編譯，以及是否出現錯誤。編譯失敗時，資料夾中可能仍保留上一次的 \code{main.pdf}，因此能開啟 PDF 並不代表新內容已經寫入。請先修正錯誤，再重新執行編譯指令。

若使用 \code{make watch}，請確認監看程序仍在執行，並且修改的檔案已經存檔。編譯成功後畫面仍未更新時，可以關閉再重新開啟 PDF，同時確認檢視器開啟的是目前專案中的 \code{main.pdf}，避免看到先前複製到其他位置的版本。若編譯結果卡在舊的狀態，執行 \code{make clean} 清掉暫存檔後再編譯一次，通常就能解決。


\end{document}
\end{lstlisting}

每一行的作用整理成\tblref{tbl:main}。

\begin{tlt}[notes={封面、審定書與三份目錄由 \texttt{\textbackslash thesisfront} 產生。}]{\texttt{main.tex} 各行的作用}{tbl:main}{L{0.5\linewidth}L{0.4\linewidth}}
內容 & 作用 \\
\midrule
\code{\documentclass{tw-thesis}} & 載入模板，讀入設定，以及圖片、表格、程式碼、假文等工具 \\
\code{\thesisfront} & 產生封面、審定書、目次、圖次、表次 \\
\code{\input{contents/front/...}} & 前置各頁，排版順序即此檔的先後順序 \\
\code{\input{contents/chapterNN}} & 正文各章，\code{\input} 的順序即為章號 \\
\code{\references} & 參考文獻，可指定文獻管理檔案位置 \\
\code{\appendix} & 此行後的章標題改用附錄編號 \code{A}、\code{B} \\
\code{\input{contents/back/...}} & 附錄 \\
\end{tlt}

幾件事需要特別注意：

\begin{enumerate}
\item 編號分界：致謝、摘要、符號列表用的是不編號的標題，所以不會變成第一章；後面用 \code{\chapter} 寫的章才會編號。
\item 章號看的是 \code{\input} 的排列順序，不是檔名，例如將 \code{chapter03} 那行移到 \code{chapter02} 前面，就會變成第二章。
\item 檔名雖然不影響結果，還是建議跟順序保持一致，免得日後自己找不到檔案。
\end{enumerate}

\section{新增與刪除章節}

模板把內容切成一章一個檔案，全部放在 \code{contents/}，檔名依章次編號，例如 \code{chapter01.tex}、\code{chapter02.tex}。改某一章不會動到其他章，編譯出錯時也可以快速判斷是哪個檔案的問題。

\code{main.tex} 裡的 \code{\input} 決定了各章的順序，新增或刪除章節時都需要自己維護這幾行。

新增一章：建立 \code{contents/chapter05.tex}（檔名依章次編號），檔案開頭放一個章標題與一個節標題，接著到 \code{main.tex} 的 \code{\references} 之前加入 \code{\input{contents/chapter05}}。

\begin{lstlisting}[language=latex]
\chapter{章標題}

\section{節標題}

內文。
\end{lstlisting}

刪除一章：把檔案刪掉，同時移除 \code{main.tex} 裡對應的那行 \code{\input}，否則 \LaTeX{} 找不到檔案會直接編譯失敗。後面各章的章號會自動遞補，正文裡的 \code{\ref} 引用也會跟著更新，都不用手改。

前置部分（\code{contents/front/}）與附錄（\code{contents/back/}）的做法相同：建立檔案，並在 \code{main.tex} 對應位置加一行 \code{\input}；刪除時也要一併移除那行。

\section{撰寫章節內容}

章節檔案建好後即可開始撰寫論文。建議先把章標題與幾個節標題列出來，再逐節填內文，不必一開始就寫完。

一章最基本的樣子如下：

\begin{lstlisting}[language=latex]
\chapter{緒論}

\section{研究背景}

第一段內文。

第二段內文，段落之間空一行。

\subsection{研究動機}

第三層標題底下的內文。
\end{lstlisting}

\code{\chapter}、\code{\section}、\code{\subsection} 依序是章、節、小節，第四層用 \code{\subsubsection}，不編號。每一層的字級、對齊方式與編號用詞都由學校設定檔決定，寫作時只要選對層級，不必自己調整外觀。章標題還會自動換頁，所以檔案開頭與結尾都不需要留空白。

段落之間空一行就會分段，首行縮排也是自動的。不要在每行末尾手動斷行，\LaTeX{} 會依版面寬度自行折行；原始碼裡的單一換行只會被當成一個空白，強制斷行反而會讓日後修改文字時出現奇怪的排版。

要讓某一章能被引用，在標題後面加上標籤，例如 \code{\chapter{研究方法}\label{cha:method}}。之後在任何地方寫 \code{\ref{cha:method}}，就會顯示成可點擊的「第三章」，章號變動時也會自動更新。圖、表與程式碼的標籤寫法相同，詳見第三章。

\section{暫時用假文看版面}

章節架構先定好、內文還沒寫時，可以先填假文，把版面、分頁與圖表位置排出來。\code{tw-thesis.sty} 包了兩個指令，預設各給五段：

\begin{lstlisting}[language=latex]
\zhlorem        % 中文五段
\enlorem        % 英文五段
\zhlorem[2]     % 顯示兩段
\enlorem[8]     % 顯示八段
\end{lstlisting}

中文假文來自 \href{https://ctan.org/pkg/zhlipsum}{zhlipsum} 套件，英文假文來自 \href{https://ctan.org/pkg/lipsum}{lipsum} 套件。兩個套件都隨 \TeX{} 發行版安裝。想換成別的文本或調整生成方式，改 \code{tw-thesis.sty} 檔尾的那兩行即可。

\chapter{版面設定}\label{cha:layout}

完成第一章的設定後，請先編譯一次，對照系所要求檢查字級、行距與頁面邊界，再調整需要修改的項目。本章會依序說明這些設定的位置與用法，建議一次修改一項，確認排版結果後再繼續。

下面的 \code{YAML} 範例都可以直接寫進 \code{config.yml} 覆蓋。若檔案裡已經有 \code{layout:}，請將新增的設定放在原有區塊下，保留範例中的縮排，不要再建立第二個同名區塊。每次修改設定後，都要重新執行 \code{python3 scripts/gen.py}（或直接用 \code{make}），\LaTeX{} 才讀得到新的值。

\section{字體與字級}

請先打開 \code{config.yml}，在 \code{layout:} 下設定本文的字型與字級。以下列出預設使用的字型；若目前顯示正常，可以先保留這些值，只調整系所要求的項目。

\begin{lstlisting}[language=yaml]
layout:
  language: "chinese" # chinese | english 決定標題語言
  font-latin: ["EB Garamond", "Tinos"] # 可填一個或多個，依序尋找
  font-cjk: ["TW-MOE-Std-Kai", "DFKai-SB"]
  font-size: 12 # pt
\end{lstlisting}

模板會依字型清單的順序尋找，用第一個找得到的字型：英文優先使用 EB Garamond，找不到才改用 Tinos（Tinos 是字寬與 Times New Roman 完全相同的開源字型），中文則優先使用\href{https://language.moe.gov.tw/material/info?m=9fe3fe82-8bbf-44c0-961d-873ea079e284}{教育部標準楷書}（授權條款為創用 CC 姓名標示─禁止改作，引用時須標示「中華民國教育部」）。這些字型檔都已經附在專案的 \code{fonts/} 裡，由 \code{tw-thesis.cls} 直接依檔案載入，不必另外安裝。兩套英文字型的授權都是 SIL OFL 1.1，可以自由散布。若系所明確要求 Times New Roman，可以把 \code{font-latin} 改成 \code{"Tinos"}：它的字寬與 Times New Roman 逐字相同，字形也相近，分頁位置不會跑掉；若一定要用原字型，因它有散布限制、模板無法附上，需先在自己的電腦安裝，再填入 \code{"Times New Roman"}。EB Garamond 的字寬與 Times New Roman 不同（平均約差 7\%），換過去後整份論文的分頁位置會改變。

若要換成系統或系所指定的其他字型，把字型名稱填進 \code{font-latin} 或 \code{font-cjk} 就好；不在 \code{fonts/} 裡的字型，需要已經安裝在系統上。這裡容易卡住的地方是：填的名稱要和字型檔內記載的家族名稱一致，而不是安裝畫面或字型網站上顯示的名稱，兩者有時不同。\tblref{tbl:font-name}列出幾種常見情境下查詢家族名稱的方式。

\begin{tlt}[notes={不確定時，一律以 \texttt{fc-list} 印出的名稱為準。}]{查詢字型家族名稱}{tbl:font-name}{L{0.28\linewidth}L{0.62\linewidth}}
情境 & 查詢方式 \\
\midrule
已經安裝在系統的字型 & 執行 \code{fc-list : family}，印出的名稱就是 \XeLaTeX{} 實際讀到的家族名，直接填入設定檔最準確。 \\
macOS 已安裝的字型 & 開啟字體簿（Font Book），點選字型後即可看到名稱。 \\
Windows 已安裝的字型 & 開啟「設定 → 個人化 → 字型」，點進字型頁面即可看到名稱。 \\
從 Google Fonts 等網站下載 & 頁面上顯示的名稱通常就是家族名，但下載後需先安裝到系統，再用 \code{fc-list} 確認一次。 \\
\end{tlt}

若要把本文改成 14 點，請在 \code{layout:} 下將 \code{font-size} 改為 \code{14}，不需要在數字後面加上單位。修改後可以找一頁中英混排的內文，檢查兩種文字的大小是否合適。章節標題有各自的字級設定，需到後面的 \code{heading:} 區塊另外調整。

\section{行高與段落}

本文的行距、段距、首行縮排與左右對齊都放在 \code{layout:} 下。可以先挑一頁有數個完整段落的內文，套用下面的設定，再比較行與行、段與段之間的距離，以及各段開頭的位置。

\begin{lstlisting}[language=yaml]
layout:
  line-height: 1.5      # 行高 1.5 倍
  par-spacing: ~        # ~ 跟隨行高；也可填 em 數值獨立調整
  first-line-indent: 2  # em，中文習慣首行縮排兩個字
  justify: true         # 左右對齊
\end{lstlisting}

調整段落時，可以先從 \code{line-height} 與 \code{first-line-indent} 開始。本文為 12 點、行高設為 \code{1.5} 時，一般文字行的基線間距為 18 點；首行縮排設為 \code{2}，則會留出兩個字寬的空間。若系所要求段落不縮排，將 \code{first-line-indent} 改為 \code{0} 即可。

\code{par-spacing} 控制兩個段落之間的間距。預設的 \code{~} 會跟隨 \code{line-height}，讓換段前後的基線距離與一般行距相同；例如 \code{line-height: 1.5} 會換算成 \code{0.5em} 的段距，加上 1em 的文字行框後，兩段之間仍維持 1.5 倍行高，不會額外多出一行空白。

若學校或系所另有規定，可以將 \code{par-spacing} 改成數字獨立覆寫，單位是 \code{em}，不需在 \code{YAML} 中加上單位；例如 \code{par-spacing: 1.2} 會加大段落間的留白。這個設定不會改變 \code{line-height}，之後調整行距時也不會跟著變動。章節原始碼中以空白行分隔的文字會形成不同段落，可以用相鄰的兩段文字觀察調整結果。

\code{config.yml} 裡還有一個 \code{line-box}，它是 Typst 版用來釘住行盒上下緣的設定；\LaTeX{} 直接以基線間距排版，用不到這一項，填了也不會有作用，保留只是為了讓兩個版本能共用同一份設定檔。

\section{邊界與紙張}

紙張大小由 \code{paper} 決定，四邊留白則在 \code{margin:} 下分別設定。以下使用 A4 紙張；若已選擇學校格式，請先查看套用結果，只加入需要覆蓋的邊界數值即可。

\begin{lstlisting}[language=yaml]
layout:
  paper: "a4"
  margin: # cm
    top: 3
    bottom: 2
    left: 3
    right: 3
\end{lstlisting}

請在 \code{margin:} 下填入四個方向的邊界，單位為公分，直接寫數字即可。例如系所要求左側留 3.5 公分，就將 \code{left} 設為 \code{3.5}，不要寫成 \code{3.5cm}。若只修改左側，其餘三邊仍會沿用學校設定或預設值，不必全部重寫。

修改後，請同時查看正文與封面，確認標題、圖表和封面資料是否仍放得下。封面的邊界由 \code{cover.margin} 另外設定，沒有填的邊會沿用正文邊界；封面各區塊另以固定間距與彈性間距安排位置。若論文題目較長，可以再調整 \code{cover:} 下的字級或間距，避免文字擠在一起。

\section{標題與編號}

章節標題的字級、對齊方式與前後間距，都集中在 \code{layout.heading} 區。下列各組清單依照標題層級排列，可以先對照章、節與小節的外觀，再調整對應位置的數值。

\begin{lstlisting}[language=yaml]
layout:
  heading:
    size: [18, 16, 14, 12]                     # pt，第 1 至 4 層
    weight: "regular"                          # regular | bold
    align: ["center", "center", "left", "left"]
    line-height: 1.2                           # 標題內部的行距
    gap-above: [1, 0.8, 0.8, 0.6]              # em
    gap-below: [1, 0.5, 0.4, 0.3]              # em
    number-gap: [1, 1, 0, 0]                   # em，編號與標題之間
\end{lstlisting}

上面的清單依序對應第一至第四層標題。例如要調整章標題的字級，就修改 \code{size} 的第一個數值；要調整節標題，則修改第二個數值。建議先保留四個位置，只更換需要的值，方便對照各層的設定。

若要修改編號前後的文字，請在 \code{config.yml} 另外加入 \code{strings:}，與 \code{layout:} 放在同一層。以下範例會讓章、節分別使用「第一章」與「第一節」的寫法：

\begin{lstlisting}[language=yaml]
strings:
  chapter: ["第", "章"]   # 第一章
  section: ["第", "節"]   # 第一節
  subsection: "、"        # 一、
  appendix: ["附錄 ", ""] # 附錄 A
\end{lstlisting}

回到章節檔案後，用 \LaTeX{} 的分節指令指定層級。\code{\chapter} 是章標題，\code{\section} 是節標題，依此類推。標題文字寫在大括號內，並依照內容的層次安排：

\begin{lstlisting}[language=latex]
\chapter{研究方法}       % 第一層，章號依出現順序自動產生
\section{實驗設計}       % 第二層，印成第一節
\subsection{資料來源}    % 第三層，印成一、
\subsubsection{前處理}   % 第四層，不編號
\end{lstlisting}

編譯時，第一層標題會自動另起一頁，章號也會依出現順序產生。因此標題只需寫「研究方法」，不必手動加上「第三章」或插入換頁指令。新增或調整章節順序後，重新編譯就會更新編號；因為目次與交叉引用要靠編號檔記錄，\LaTeX{} 需要連續編譯兩次才會完全正確，\code{latexmk} 會自動處理。

\section{圖表與程式碼的標題位置}\label{sec:caption}

圖片、表格與程式碼的標題位置可以分別設定。請在 \code{config.yml} 的 \code{layout:} 下找到 \code{caption-position:}，將要調整的項目設為 \code{"top"} 或 \code{"bottom"}，分別代表放在內容上方或下方。以下是目前的預設值：

\begin{lstlisting}[language=yaml]
layout:
  caption-position:
    image: "bottom"    # 圖片標題放在下方
    table: "bottom"    # 表格標題放在下方
    code: "bottom"     # 程式碼標題放在下方
\end{lstlisting}

例如要將所有表格的標題移到上方，只需把 \code{table} 改為 \code{"top"}，圖片與程式碼不會跟著改變。圖片的設定也會套用到並排與網格圖片的總標題。修改後重新編譯，檢查標題與內容的位置即可，不需要逐一修改各章的圖表。

\section{目次與頁碼}

目次的收錄層級、縮排與頁碼格式都放在 \code{outline:} 下，與 \code{layout:} 位於同一層。以下列出常用設定，初次使用時可以先保留預設值，等章節內容齊全後再調整。

\begin{lstlisting}[language=yaml]
outline:
  depth: 3            # 目次收到第幾層
  indent: 1           # em，每層縮排一個字
  title-size: 16      # pt，目次、圖次、表次的字級
  page-numbering: "i" # 前置部分：i, ii, iii
  body-numbering: "1" # 正文開始：1, 2, 3
  entry-gap: 1        # em，章與章之間的間距
  entry-bold: true    # 章標題在目次裡加粗
\end{lstlisting}

目次會依照章節標題自動產生，圖次與表次則整理文件中的圖片和表格。若只想在目次中列出章與節，將 \code{depth} 改為 \code{2} 即可，正文的標題仍會保留原有層級。編譯後可以點擊目次條目，確認是否跳到對應位置。

要修改這三頁的標題，請在 \code{titles:} 下設定文字。字串中的 \code{~} 會轉成字距，例如 \code{"目~次"} 會拉開兩個字的間隔；若要直接顯示「目錄」，填入 \code{"目錄"} 即可：

\begin{lstlisting}[language=yaml]
titles:
  toc: "目~次"   # 有些學校用「目錄」
  lof: "圖~次"
  lot: "表~次"
\end{lstlisting}

頁碼的位置與顯示文字放在 \code{page-number:} 下，與 \code{outline:} 位於同一層。以下都是預設值，不修改也能直接使用：

\begin{lstlisting}[language=yaml]
page-number:
  position: "bottom" # top：頁首；bottom：頁尾
  align: "center"    # left：靠左；center：置中；right：靠右
  format: "{page}"   # {page} 會替換成該頁的頁碼
\end{lstlisting}

各項設定可依下表調整。將 \code{position} 設為 \code{"top"}、\code{align} 設為 \code{"right"}，頁碼就會移到頁首靠右。

\begin{tlt}{頁碼設定參數}{tbl:page-number}{L{0.15\linewidth}L{0.75\linewidth}}
參數 & 用途與範例 \\
\midrule
\code{position} & 位置：\code{"top"} 頁首、\code{"bottom"} 頁尾（預設）。 \\
\code{align} & 對齊：\code{"left"}、\code{"center"}（預設）、\code{"right"}。 \\
\code{format} & 顯示文字：\code{"第 {page} 頁"}；設為 \code{""} 可隱藏頁碼。 \\
\end{tlt}

\code{format} 中的 \code{{page}} 會替換成頁碼，例如 \code{"第 {page} 頁"} 顯示為「第 1 頁」；也可以寫成 \code{"論文 {page}"}。其他文字會直接顯示，特殊符號（如 \code{%}、\code{&}）由產生設定的程式自動處理，不需要加反斜線。

\code{outline.page-numbering} 與 \code{outline.body-numbering} 分別控制前置部分與正文的編號方式，預設為小寫羅馬數字與阿拉伯數字。\code{page-number.format} 只調整頁首或頁尾的文字，目次仍保留原本的頁碼格式。封面與審定書不顯示頁碼。完成章節調整後，請檢查正文第一頁與目次所列的頁碼，確認兩者對應正確。

\section{浮水印與定稿}

提交前若需要加上浮水印，請先檢查學校設定是否已指定圖片。需要自行指定時，可以在 \code{config.yml} 加入下列區塊，將 \code{image} 改成要使用的圖檔路徑，再調整大小與深淺。

\begin{lstlisting}[language=yaml]
watermark:
  image: "/assets/watermarks/nccu.pdf"
  width: 100    # %，佔頁寬的比例，置中
  opacity: 100  # %，100 是原圖深淺，50 會刷淡一半
\end{lstlisting}

設定好圖片後，執行 \code{make release}，再打開 PDF 的正文頁面查看浮水印。一般的 \code{make} 不會啟用浮水印，封面與前置頁面也不會加上。若 \code{image} 的值是 \code{~}，表示沒有指定圖片，即使執行 \code{make release} 也不會顯示。

若要換成自己的浮水印，請先將圖片放進 \code{assets/watermarks/}，再修改 \code{image} 的路徑。可以用 \code{width} 調整圖片占頁寬的比例，或降低 \code{opacity}，讓浮水印淡一些。是否需要自行加上浮水印，請依提交時的學校要求處理。

沒有安裝 \code{make} 時，也可以在專案根目錄執行以下指令，輸出的檔案同樣是 \code{main.pdf}：

\begin{lstlisting}[language=bash]
python3 scripts/gen.py
latexmk -g -xelatex -usepretex='\def\twwatermark{1}' main.tex
\end{lstlisting}


% references
\references

% appendices
\appendix
\chapter{編譯問題排查}\label{appendix-a}

使用模板時，若遇到編譯失敗、圖片找不到或 PDF 沒有更新，可以先依本附錄的順序檢查。建議保留終端機中的錯誤訊息，從第一個錯誤指出的檔案與位置開始處理，修正後再重新編譯。

這份檔案也示範附錄的排版方式。換成自己的論文時，可以保留章標題的寫法，將標題與內文替換成問卷、補充數據或其他資料，再到 \code{main.tex} 確認附錄的裝訂順序。

\section{指令執行失敗}

若終端機顯示找不到 \code{xelatex} 或 \code{latexmk}，請先確認已完成安裝，再重新開啟終端機，執行

\begin{lstlisting}[language=bash]
xelatex --version
\end{lstlisting}

確認能夠顯示版本資訊後，再嘗試編譯。若找不到的是 \code{make}，則可以先使用以下指令，直接產生 PDF：

\begin{lstlisting}[language=bash]
texlua scripts/gen.lua
latexmk -xelatex main.tex
\end{lstlisting}

若錯誤訊息指出找不到 \code{main.tex} 或 \code{.tw-thesis-config.tex}，請檢查終端機目前所在的資料夾。編譯指令應在專案根目錄執行，也就是能看到 \code{main.tex}、\code{config.yml} 與 \code{contents/} 的位置；\code{.tw-thesis-config.tex} 則是 \code{scripts/gen.lua} 產生的，先執行過一次就會出現。若目前位於章節資料夾，請先切換回根目錄再試一次。

\section{修改設定後無法編譯}

修改設定後出現錯誤時，請先看訊息指出哪個檔案、哪一個欄位。若包含 \code{unknown key}，可以對照 \code{schools/default.yml} 檢查欄位拼字與所在層級；若是 \code{YAML} 解析錯誤，則先確認縮排使用空白，且冒號後保留空白，再檢查引號是否成對。

加入新設定時，請先找出檔案中是否已有相同區塊，避免在同一層重複建立 \code{layout:}。例如要同時修改字級與左側邊界，應把兩項設定放在同一個 \code{layout:} 下，並讓 \code{left} 再往內縮排一層：

\begin{lstlisting}[language=yaml]
layout:
  font-size: 12
  margin:
    left: 3.5
\end{lstlisting}

若一次修改了多個設定，暫時無法判斷錯誤來源，可以先將剛加入的內容複製到別處保存，再移除這些設定，確認原本的文件能正常編譯。之後逐項加回，每次存檔後檢查結果，較容易找到需要修正的位置。

\section{圖片或引用找不到}

圖片無法載入時，請先確認檔案確實放在 \code{figures/} 的指定位置，再比對路徑中的檔名、副檔名與大小寫。路徑是從專案根目錄算起的，不是從目前章節所在的資料夾開始。若曾經移動或重新命名圖片，也要一併修改章節中的路徑。

圖表引用顯示成 \code{??} 時，請檢查圖表是否已加上 \code{\label}，以及名稱是否與內文的 \code{\ref} 完全相同；也別忘了 \LaTeX{} 需要編譯兩次才會解出交叉引用，\code{latexmk} 會自動處理。若是文獻引用出錯，則到 \code{.bib} 檔案確認識別名稱，再檢查 \code{config.yml} 的 \code{bibliography.path} 是否包含該文獻檔。修正後重新編譯，確認引用已顯示為編號或文獻資料。

\section{PDF 沒有更新}

修改內容後，請先確認終端機是否已完成這次編譯，以及是否出現錯誤。編譯失敗時，資料夾中可能仍保留上一次的 \code{main.pdf}，因此能開啟 PDF 並不代表新內容已經寫入。請先修正錯誤，再重新執行編譯指令。

若使用 \code{make watch}，請確認監看程序仍在執行，並且修改的檔案已經存檔。編譯成功後畫面仍未更新時，可以關閉再重新開啟 PDF，同時確認檢視器開啟的是目前專案中的 \code{main.pdf}，避免看到先前複製到其他位置的版本。若編譯結果卡在舊的狀態，執行 \code{make clean} 清掉暫存檔後再編譯一次，通常就能解決。


\end{document}
\end{lstlisting}

What each line does is summarized in \tblref{tbl:main}.

\begin{tlt}[notes={The cover, verification letter and the three lists are produced by \texttt{\textbackslash thesisfront}.}]{What each line of \texttt{main.tex} does}{tbl:main}{L{0.5\linewidth}L{0.4\linewidth}}
Content & Effect \\
\midrule
\code{\documentclass{tw-thesis}} & Loads the template, reads the settings, and loads the tools: figures, tables, code, filler text \\
\code{\thesisfront} & Produces the cover, verification letter, table of contents, list of figures and list of tables \\
\code{\input{contents/front/...}} & The front-matter pages, in the order of this file \\
\code{\input{contents/chapterNN}} & The chapters; the order of the \code{\input} lines is the chapter number \\
\code{\references} & The references; the bibliography files are set in the config \\
\code{\appendix} & Chapters after this line switch to appendix numbering \code{A}, \code{B} \\
\code{\input{contents/back/...}} & The appendices \\
\end{tlt}

A few things need special attention:

\begin{enumerate}
\item Where numbering starts: the acknowledgements, abstract and notation list use unnumbered headings, so they do not become Chapter 1; chapters written with \code{\chapter} after them are numbered.
\item The chapter number follows the order of the \code{\input} lines, not the file name. For example, moving the \code{chapter03} line above \code{chapter02} makes it the second chapter.
\item File names do not affect the result, but keeping them in line with the order saves you from losing track of files later.
\end{enumerate}

\section{Adding and Removing Chapters}

The template splits the content into one file per chapter, all inside \code{contents/}, named by chapter number, for example \code{chapter01.tex} and \code{chapter02.tex}. Editing one chapter does not touch the others, and when compiling fails you can quickly tell which file is at fault.

The \code{\input} lines in \code{main.tex} decide the chapter order, so you maintain those lines yourself when adding or removing chapters.

To add a chapter: create \code{contents/chapter05.tex} (named by chapter number), put a chapter heading and a section heading at the top, then add \code{\input{contents/chapter05}} in \code{main.tex} before \code{\references}.

\begin{lstlisting}[language=latex]
\chapter{Chapter Title}

\section{Section Title}

Body text.
\end{lstlisting}

To remove a chapter: delete the file and remove the matching \code{\input} line from \code{main.tex}, otherwise \LaTeX{} cannot find the file and compiling fails. The chapters after it are renumbered automatically, and the \code{\ref} references in the text update too; none of that needs manual edits.

Front matter (\code{contents/front/}) and appendices (\code{contents/back/}) work the same way: create the file and add an \code{\input} line at the right place in \code{main.tex}; when deleting, remove that line as well.

\section{Writing Chapter Content}

Once a chapter file exists, you can start writing. It is a good idea to list the chapter title and a few section headings first and then fill in the text section by section; there is no need to write everything at once.

The most basic chapter looks like this:

\begin{lstlisting}[language=latex]
\chapter{Introduction}

\section{Background}

The first paragraph.

The second paragraph; leave a blank line between paragraphs.

\subsection{Motivation}

Text under a third-level heading.
\end{lstlisting}

\code{\chapter}, \code{\section} and \code{\subsection} are chapter, section and subsection in that order; the fourth level uses \code{\subsubsection} and is not numbered. The font size, alignment and numbering words of each level come from the school settings, so when writing you only choose the right level and never adjust the look yourself. A chapter heading also starts a new page automatically, so leave no blank space at the beginning or end of the file.

A blank line starts a new paragraph, and the first-line indent is automatic. Do not break lines by hand at the end of each line: \LaTeX{} wraps according to the page width, and a single line break in the source counts as one space, so forcing breaks only causes odd layouts when you edit the text later.

To make a chapter referable, add a label after its heading, for example \code{\chapter{Method}\label{cha:method}}. Writing \code{\ref{cha:method}} anywhere then shows a clickable ``Chapter 3'', which updates when the chapter number changes. Labels for figures, tables and code work the same way; see Chapter 3.

\section{Filler Text for Trying the Layout}

Once the chapter structure is settled but the text is not written, you can fill it with filler text to lay out the pages, page breaks and figure and table positions. \code{tw-thesis.sty} wraps two commands, five paragraphs each by default:

\begin{lstlisting}[language=latex]
\zhlorem        % five paragraphs of Chinese
\enlorem        % five paragraphs of English
\zhlorem[2]     % two paragraphs
\enlorem[8]     % eight paragraphs
\end{lstlisting}

The Chinese filler comes from the \href{https://ctan.org/pkg/zhlipsum}{zhlipsum} package and the English filler from \href{https://ctan.org/pkg/lipsum}{lipsum}; both ship with the \TeX{} distribution. To use different text or change how it is generated, edit those two lines at the end of \code{tw-thesis.sty}.
